\documentclass[12pt,a4paper]{report}
\usepackage[utf8]{inputenc}
\usepackage[T1]{fontenc}
\usepackage[french]{babel}
\usepackage{geometry}
\geometry{margin=2.5cm}
\usepackage{setspace}
\usepackage{graphicx}
\usepackage{float}
\usepackage{array}
\usepackage{xltabular}
\usepackage{tabularx}
\usepackage{xcolor}
\definecolor{bleu}{rgb}{0,0,1}
\usepackage{url}
\setlength{\parskip}{1em}
\setlength{\parindent}{0em}
\usepackage{amsmath}
\usepackage{amssymb}
\onehalfspacing
\usepackage{tikz}
\usetikzlibrary{shapes.geometric, arrows, positioning, arrows.meta}
\usepackage{booktabs}
\usepackage{algorithm}
\usepackage{algpseudocode}
\usepackage[normalem]{ulem}
\usepackage[pdfborder={0 0 1}]{hyperref}
\usepackage{subcaption}

\newcommand{\encours}[1]{\textcolor{red}{\textbf{[EN COURS -- #1]}}}
\newcommand{\todo}[1]{\textcolor{orange}{\textit{[À COMPLÉTER : #1]}}}

\begin{document}

\begin{titlepage}
\vspace*{0.95cm}
\begin{center}
\textnormal{\textbf{Université de Mons}}\\[0.3em]
\textnormal{\textbf{Faculté des Sciences}}\\[0.3em]
\textnormal{\textbf{Département d'Informatique}}\\[0.3em]
\textnormal{\textbf{Finalité spécialisée en Artificial Intelligence and Data Analytics}}\\[0.3em]
\end{center}
\vspace*{3cm}
\centering

{\Large \textbf{Apprentissage profond et fédéré pour la classification de la
densité mammaire pour la prédiction du cancer du sein.}}\\[0.3cm]
{\large Construire un modèle capable de classer automatiquement la densité
mammaire selon BI-RADS à partir d'images mammographiques}\\[2cm]

\begin{tabular}{@{} l l @{}}
\textbf{Étudiant :} & ALLA Abdelouahad \\
\textbf{Maître de mémoire :} & Xavier Lessage, \textbf{CETIC} \\
\textbf{Directeur de mémoire :} & MAHMOUDI Saïd, \textbf{UMONS} \\
\end{tabular}

\textit{Année académique 2025--2026}

\vfill

\noindent
\includegraphics[height=1.8cm]{image/UMONS-Logo1.jpg}
\hfill
\includegraphics[height=1.8cm]{image/FS-Logo.jpg}
\hfill
\includegraphics[height=1.8cm]{image/FPMS-Logo.jpg}

\end{titlepage}

\chapter*{Remerciements}
\addcontentsline{toc}{chapter}{Remerciements}

Je tiens à exprimer ma profonde gratitude à mon maître de mémoire, M.
Xavier Lessage (CETIC), pour son suivi régulier tout au long de ce travail.
Son orientation, son soutien et ses encouragements ont été cruciaux pour
l'aboutissement de ce travail, et nos échanges réguliers ont permis de
garder une cohérence scientifique tout au long des expérimentations, en
orientant certains choix méthodologiques abordés dans ce mémoire.

Je tiens également à adresser mes sincères remerciements à mon directeur
de mémoire, le professeur Saïd Mahmoudi (UMONS), pour ses retours et sa
disponibilité.

\tableofcontents

% ============================================================================
\chapter*{Introduction}
\addcontentsline{toc}{chapter}{Introduction}

Le cancer du sein touche encore aujourd'hui énormément de femmes à travers
le monde, et reste un problème de santé publique majeur. Les traitements se
sont beaucoup améliorés ces dernières décennies. Mais ce qui pèse le plus
sur les chances de guérison reste le moment du diagnostic : plus une tumeur
est repérée tôt, plus les traitements ont des chances d'être efficaces.

Le dépistage repose aujourd'hui surtout sur la mammographie. Cet examen
n'offre pourtant pas la même fiabilité selon les patientes, à cause de la
composition du tissu mammaire. Chez certaines femmes, le sein contient
beaucoup de tissu fibreux et glandulaire -- on parle alors de sein
\emph{dense} -- qui apparaît en blanc sur le cliché, exactement comme les
lésions suspectes. Cette ressemblance visuelle peut littéralement cacher une
tumeur au milieu d'un tissu sain, et cela touche justement les patientes les
plus à risque : une densité élevée augmente elle-même les chances de
développer un cancer.

Pour limiter cette part de subjectivité, le système BI-RADS classe la
densité mammaire en quatre catégories, d'un sein presque entièrement
graisseux à un sein extrêmement dense. Ce classement reste néanmoins établi
à l'œil par le radiologue. Il existe donc une variabilité d'un praticien à
l'autre, et parfois même d'un examen à l'autre pour la même patiente.

C'est cette limite que ce mémoire cherche à dépasser, en s'appuyant sur
l'apprentissage profond pour estimer la densité mammaire de façon
automatique et reproductible, sans dépendre du jugement d'un seul
observateur. Le sujet impose aussi d'explorer une seconde dimension, liée
aux contraintes réelles du monde hospitalier : la confidentialité des
données médicales. C'est le rôle de l'apprentissage fédéré, qui permet en
principe d'entraîner un modèle sur plusieurs sources de données sans jamais
centraliser les images elles-mêmes.

Conformément au titre de ce mémoire, \emph{Apprentissage profond et fédéré
pour la classification de la densité mammaire pour la prédiction du cancer
du sein}, le travail se structure en deux grands axes. Le premier, présenté
aux Chapitres~\ref{sec:etat-art-positionnement} et~\ref{ch:approches},
regroupe les architectures d'apprentissage profond (Vision Transformer,
CNN, fusions multi-branches) entraînées de façon classique et centralisée
sur VinDr-Mammo. Le second, présenté au Chapitre~\ref{sec:federe}, explore
l'apprentissage fédéré (FedAvg) entre deux sources de données distinctes
(VinDr-Mammo et DDSM/CBIS-DDSM), sans jamais centraliser les images
elles-mêmes -- une contrainte réaliste pour des données médicales, où le
partage direct entre institutions se heurte à des obligations de
confidentialité.

Le chapitre suivant situe d'abord ce travail par rapport à la littérature
existante (Chapitre~\ref{sec:etat-art-positionnement}). Le détail de chaque
approche testée, avec ses résultats propres, vient ensuite
(Chapitre~\ref{ch:approches}), suivi de l'apprentissage fédéré
(Chapitre~\ref{sec:federe}) et de l'analyse d'interprétabilité du modèle
retenu.

% ============================================================================
\chapter{Contexte médical et classification de la densité mammaire}
\label{ch:contexte-medical}

\section{Rappel sur le cancer du sein}

Le cancer du sein correspond à une prolifération incontrôlée de cellules qui
se trouvaient initialement dans la glande mammaire. Sans traitement, cette
prolifération peut dépasser son site d'origine, envahir les tissus voisins,
puis atteindre des organes distants par voie sanguine ou lymphatique, sous
forme de métastases. On distingue classiquement les tumeurs
\textbf{bénignes}, qui restent localisées et sans danger immédiat, des
tumeurs \textbf{malignes}, capables de cette évolution invasive.

\subsection{Organisation générale du sein}

Le sein est un organe pair, situé à l'avant du thorax et posé sur le muscle
pectoral qui lui sert de support~\cite{ref5}. Il associe trois types de
tissus : la glande mammaire elle-même, un stroma de soutien fait de tissu
conjonctif, et une composante graisseuse dont la proportion varie fortement
d'une femme à l'autre. Le tout est recouvert par la peau.

Sur le plan fonctionnel, la glande se structure en 15 à 20 \textbf{lobes},
eux-mêmes divisés en 20 à 40 \textbf{lobules} : ce sont les unités qui
produisent le lait pendant l'allaitement. Chaque lobule se draine dans un
réseau de \textbf{canaux galactophores}, qui achemine le lait vers le
mamelon. C'est justement le long de ce réseau que naissent le plus souvent
les cancers du sein, d'où la fréquence des carcinomes dits canalaires par
rapport aux formes lobulaires.

La Figure~\ref{fig:anatomie_sein} reprend ces structures sous forme de
coupe.

\begin{figure}[H]
    \centering
    \includegraphics[width=0.70\textwidth]{image/anatomie_sein_.png}
    \caption{Coupe schématique du sein : (\textcolor{blue}{1})~paroi
    thoracique, (\textcolor{blue}{2})~muscle pectoral,
    (\textcolor{blue}{3})~lobules, (\textcolor{blue}{4})~lobes,
    (\textcolor{blue}{5})~mamelon, (\textcolor{blue}{6})~aréole,
    (\textcolor{blue}{7})~canal galactophore, (\textcolor{blue}{8})~tissu
    adipeux, (\textcolor{blue}{9})~tissu conjonctif~\cite{ref16}.}
    \label{fig:anatomie_sein}
\end{figure}

\subsection{Origine et types de tumeurs mammaires}

Un cancer du sein se développe presque toujours à partir de l'un de ces deux
compartiments : les lobules ou les canaux galactophores~\cite{ref6}. Quand
les cellules anormales restent confinées à leur tissu d'origine, on parle de
tumeur \textbf{in situ} (la membrane basale n'est pas franchie). Quand elles
envahissent le tissu environnant, avec un risque de dissémination à
distance, on parle de tumeur \textbf{invasive}. En croisant cette
distinction avec le site d'origine, on obtient les quatre présentations
cliniques usuelles : carcinome canalaire in situ, carcinome canalaire
invasif, carcinome lobulaire in situ et carcinome lobulaire invasif. Cette
classification conditionne directement le pronostic et le choix
thérapeutique, comme le résume la Figure~\ref{fig:stades_cancer}.

\begin{figure}[H]
    \centering
    \includegraphics[width=0.90\textwidth]{image/stades_cancer.png}
    \caption{Les quatre présentations principales du cancer du sein selon
    leur localisation d'origine (canalaire ou lobulaire) et leur caractère
    in situ ou invasif~\cite{ref16}.}
    \label{fig:stades_cancer}
\end{figure}

\section{La densité mammaire}

\subsection{Ce que mesure la densité}

Contrairement à ce que son nom pourrait laisser penser, la densité mammaire
ne décrit pas une sensation au toucher, mais un contraste visuel observé sur
la mammographie~\cite{ref2}. Le tissu adipeux, radiotransparent, apparaît
sombre sur le cliché. Le tissu fibreux et glandulaire, lui, est
radio-opaque et apparaît clair~\cite{ref1}. La densité mammaire correspond
donc à la proportion de ce second type de tissu par rapport au premier.

Cette proportion compte pour deux raisons. D'une part, une lésion cancéreuse
apparaît elle aussi en blanc sur la mammographie : plus le tissu environnant
est dense, plus il devient difficile de distinguer une anomalie réelle du
bruit visuel ambiant. C'est cet effet de masquage qui réduit la sensibilité
du dépistage~\cite{ref3}. D'autre part, une densité élevée est associée à un
risque accru de cancer du sein, indépendamment de cet effet de masquage --
même si les mécanismes biologiques précis restent débattus dans la
littérature~\cite{ref2, ref4}.

\subsection{La classification BI-RADS}

Pour que cette évaluation ne dépende pas entièrement du jugement individuel
de chaque radiologue, l'American College of Radiology a défini un cadre de
référence, le \emph{Breast Imaging Reporting and Data System} (BI-RADS), qui
répartit la densité mammaire en quatre catégories~\cite{ref16} :

\begin{itemize}
    \item \textbf{Catégorie A (ACR1) -- sein presque entièrement adipeux :}
    la composante fibroglandulaire ne dépasse pas 25\,\%. Le contraste
    radiologique est excellent, ce qui rend la détection d'anomalies
    particulièrement fiable.
    \item \textbf{Catégorie B (ACR2) -- densité fibroglandulaire éparse :}
    entre 25\,\% et 50\,\% de tissu dense, réparti de façon discontinue sur
    un fond globalement clair. La sensibilité de l'examen reste bonne.
    \item \textbf{Catégorie C (ACR3) -- sein hétérogène dense :} entre
    50\,\% et 75\,\% de tissu dense. L'interprétation devient plus délicate
    car certaines petites lésions peuvent se confondre avec le tissu
    environnant.
    \item \textbf{Catégorie D (ACR4) -- sein extrêmement dense :} plus de
    75\,\% de tissu dense. C'est la configuration la plus difficile à lire
    pour le radiologue, et paradoxalement celle où le risque intrinsèque de
    cancer est le plus élevé.
\end{itemize}

\begin{figure}[H]
    \centering
    \begin{subfigure}[b]{0.22\textwidth}
        \centering
        \includegraphics[width=\textwidth]{image/acr1-a.png}
        \caption*{ACR 1 (A)}
    \end{subfigure}
    \hfill
    \begin{subfigure}[b]{0.22\textwidth}
        \centering
        \includegraphics[width=\textwidth]{image/acr2-b.png}
        \caption*{ACR 2 (B)}
    \end{subfigure}
    \hfill
    \begin{subfigure}[b]{0.22\textwidth}
        \centering
        \includegraphics[width=\textwidth]{image/acr3-c.png}
        \caption*{ACR 3 (C)}
    \end{subfigure}
    \hfill
    \begin{subfigure}[b]{0.22\textwidth}
        \centering
        \includegraphics[width=\textwidth]{image/acr4-d.png}
        \caption*{ACR 4 (D)}
    \end{subfigure}
    \caption{Exemple des quatre catégories de densité BI-RADS, illustré à
    partir d'images réelles du jeu de données VinDr-Mammo utilisé dans ce
    mémoire.}
    \label{fig:birads_categories}
\end{figure}

Cette dépendance à l'interprétation visuelle du radiologue, et la
variabilité inter- et intra-observateur qu'elle entraîne, est précisément ce
que ce mémoire propose de dépasser. L'idée est de s'appuyer sur
l'apprentissage profond pour produire une estimation automatique,
reproductible et indépendante de l'observateur.

% ============================================================================
\chapter{Positionnement par rapport à l'état de l'art}
\label{sec:etat-art-positionnement}

Définition et enjeux. Le sein est composé de trois types de tissus :
adipeux, glandulaire et fibreux ; la proportion de ces tissus détermine la
densité mammaire. Les tissus denses (glandulaires et fibreux) apparaissent
en blanc sur les mammographies, tandis que les tissus adipeux, plus
sombres, traduisent une faible densité. Une densité élevée constitue un
enjeu clinique majeur : elle peut masquer des lésions (difficulté de
diagnostic) et constitue un facteur de risque reconnu du cancer du sein.
Au-delà de la classification en elle-même, cet enjeu s'inscrit directement
dans l'objectif plus large de ce mémoire : une estimation fiable et
automatique de la densité contribue à la prédiction du risque de cancer du
sein, la densité étant elle-même un facteur de risque indépendant reconnu
dans la littérature~\cite{ref2, ref4}. C'est cette double perspective --
classification et évaluation du risque -- qui structure la revue de
littérature qui suit.

\section{Apprentissage fédéré en imagerie médicale}
\label{sec:etat-art-federe}

Le second axe de ce mémoire -- l'apprentissage fédéré -- répond à une
contrainte spécifique au monde médical que les méthodes de classification
présentées dans la suite de ce chapitre ne traitent pas : la plupart des
jeux de données utilisés dans la littérature (y compris VinDr-Mammo dans ce
travail) proviennent d'un nombre restreint de cohortes déjà centralisées,
alors qu'en pratique les données médicales réelles restent réparties entre
institutions, pour des raisons légales et éthiques de confidentialité
\cite{ref_rieke_fl_digital_health}. L'algorithme FedAvg de McMahan et al.\
\cite{ref_fedavg}, déjà présenté en détail au Chapitre~\ref{sec:federe},
propose une réponse technique à ce problème : entraîner un modèle
collaboratif sans jamais centraliser les images elles-mêmes, seuls les
poids des modèles locaux étant échangés.

Plusieurs travaux ont depuis appliqué ce principe à l'imagerie médicale.
Sheller et al.\ \cite{ref_sheller_fl_brain} démontrent la faisabilité d'un
entraînement fédéré multi-institutionnel pour la segmentation de tumeurs
cérébrales, sans qu'aucun hôpital participant n'ait besoin de partager ses
données brutes avec les autres. Kaissis et al.\
\cite{ref_kaissis_fl_privacy} passent en revue les enjeux de sécurité et de
confidentialité propres à l'apprentissage fédéré appliqué à l'image
médicale, un point directement pertinent pour la motivation de ce mémoire.
Plus proche encore de notre sujet précis, Roth et al.\
\cite{ref_roth_fl_breast_density} rapportent une implémentation réelle
d'apprentissage fédéré pour la classification de la densité mammaire entre
plusieurs institutions hospitalières -- c'est le précédent le plus direct
pour l'expérimentation FedAvg menée dans ce mémoire (Chapitre~\ref{sec:federe}),
à la différence près que leur déploiement s'appuie sur de vrais serveurs
distincts par institution, alors que notre implémentation reste une
simulation sur une seule machine (voir la précision méthodologique de la
Section~\ref{sec:federe}).

\section{Méthodes classiquee}

Avant que le Deep Learning ne s'impose, la détermination de la densité
mammaire s'appuyait surtout sur des descripteurs statistiques et texturaux,
combinés à des algorithmes de classification supervisée. On peut notamment
citer :\\
\begin{itemize}
    \item \textbf{L'utilisation de réseaux de neurones artificiels appliqués aux histogrammes :} cette technique classe directement le niveau de densité à partir de la répartition des niveaux de gris, sans étape de segmentation. Elle atteint un taux de précision avoisinant les 87\,\%.
\item \textbf{La segmentation par Fuzzy C-Means associée à une combinaison bayésienne :} cette méthode distingue les régions denses des zones adipeuses, avant d'appliquer des modèles de classification hybrides (kNN, arbres de décision), dont les résultats sont comparés aux jugements d'experts radiologues.
\item \textbf{La segmentation topographique couplée à l'algorithme kNN :} cette approche repose sur une représentation morphologique hiérarchisée du sein, ce qui permet de repérer les zones à forte densité avec une bonne fiabilité.
\end{itemize}
Ces techniques ont l'avantage d'être simples et interprétables. Elles
souffrent cependant de plusieurs faiblesses : une forte dépendance à la
qualité des images, une difficulté à représenter les relations spatiales
complexes, et des performances qui s'essoufflent sur de grands jeux de
données. C'est pour ces raisons qu'elles ont progressivement cédé la place
aux méthodes d'apprentissage profond, capables d'extraire seules des
représentations hiérarchiques directement à partir des données brutes.

\section{Méthodes par Deep Learning et hybrides}
Les architectures profondes, en particulier les CNN et les Transformers,
ont changé la façon d'aborder la classification de la densité mammaire.

\subsection{CNN et approches hiérarchiques}

Plusieurs travaux organisent la classification en étapes successives plutôt
qu'en un seul classifieur direct. Trace4BDensity \cite{trace4bdensity}
enchaîne trois réseaux ResNet-50 en cascade -- d'abord tissu dense A/B
contre C/D, puis A contre B, puis C contre D -- pour atteindre 89\,\%
d'exactitude. Kumar et al. \cite{ref22} combinent de leur côté un ensemble
de classifieurs neuronaux hiérarchiques et des descripteurs de texture GLCM,
ce qui fait progresser le taux de réussite de 79,5\,\% à 90,8\,\%. Cette
logique de décision par étapes, plus proche du raisonnement clinique par
élimination, est directement celle que nous reprenons avec notre
architecture de classification hiérarchique en cascade (Approche 7, voir
Section~\ref{sec:correspondance}).

\subsection{Segmentation et modèles hybrides}

D'autres approches passent par une étape de segmentation explicite avant la
classification. Grâce à leurs convolutions dilatées et à leurs mécanismes
multi-échelles, les architectures DeepLabV3 et U-Net permettent de délimiter
plus précisément les régions denses du tissu mammaire. La méthode par
segmentation cGAN de Ramirez et al. \cite{ref18} va plus loin en chaînant
segmentation et classification dans un seul pipeline end-to-end, ce qui
améliore à la fois la précision (98\,\%) et l'explicabilité du résultat.

\subsection{Transformers et architectures récentes}

L'arrivée des Transformers en vision (ViT, Swin) a introduit un mécanisme
d'attention capable de capturer des relations globales dans l'image, en
complément des filtres convolutifs qui restent, eux, des détecteurs
essentiellement locaux :
\begin{itemize}
    \item \textbf{MammoViT} \cite{ref19} associe ResNet50 et ViT et atteint
    97,4\,\% en combinant caractéristiques locales et contexte global --
    l'architecture qui a directement inspiré notre Approche 2 (fusion
    ViT/ResNet50, voir Figure~\ref{fig:mammovit_architecture}) ;
    \item le \textbf{Swin Transformer} \cite{ref20}, à attention
    hiérarchique, atteint 74,9\,\% sur une classification à 4 classes
    BI-RADS -- un score nettement plus proche des nôtres que les résultats
    binaires évoqués plus haut, ce qui confirme que l'écart vient surtout du
    passage au multi-classe (voir Section~\ref{sec:ecart-litterature}
    ci-dessous) ;
    \item des modèles multi-vues comme MaMVT exploitent conjointement les
    incidences CC et MLO via une attention croisée entre les deux vues, une
    idée que l'on retrouve -- sous une forme plus simple, par fusion de
    caractéristiques plutôt que par attention croisée -- dans notre
    Approche 6.
    
\end{itemize}
\begin{figure}[H]
    \centering
     \includegraphics[width=0.75\textwidth]{image/MammoViT_model.png}
    \caption{Architecture MammoViT, reprise de l'état de l'art -- référence
    directe pour notre Approche 2.}
    \label{fig:mammovit_architecture}
\end{figure}
\subsection{Écart avec les meilleurs résultats publiés}
\label{sec:ecart-litterature}

Nos meilleurs résultats se situent actuellement entre 73\,\% et 87\,\%
selon l'approche. Plusieurs travaux de la littérature annoncent pourtant des
scores nettement supérieurs : MammoViT \cite{ref19} rapporte 97,4\,\% de
précision grâce à la fusion des caractéristiques locales et globales,
l'approche par segmentation cGAN de Ramirez et al. \cite{ref18} annonce
98\,\%, et l'ensemble de classifieurs neuronaux de Kumar et al.
\cite{ref22} atteint 90,8\,\%. Plutôt que de voir cet écart comme un échec,
il vaut mieux le décomposer :
\begin{itemize}
    \item \textbf{Le jeu de données diffère} : contrairement à nous, une
    partie de ces travaux s'appuient sur des cohortes cliniques privées déjà
    nettoyées et annotées par un consensus de radiologues (Lehman et al.
    \cite{ref13}, Ho et al. \cite{ref3}), ou sur des versions retravaillées
    de DDSM. Nous entraînons directement sur VinDr-Mammo brut, sans
    curation manuelle préalable, ce qui expose davantage le modèle au bruit
    d'annotation et au déséquilibre naturel des classes.
    \item \textbf{La tâche n'est pas toujours identique} : plusieurs de ces
    scores élevés concernent une classification \emph{binaire}
    (dense / non dense), une tâche structurellement plus facile que notre
    classification à 4 classes ACR. Le tableau comparatif
    (Table~\ref{tab:comparatif_deep_learning} de l'état de l'art) le montre
    bien : dès qu'on repasse en multi-classe, les scores se rapprochent
    davantage des nôtres.
    \item \textbf{L'architecture exacte n'est pas toujours reproductible} :
    les hyperparamètres précis, le prétraitement, et parfois l'architecture
    elle-même ne sont pas entièrement décrits dans les publications.
\end{itemize}

\subsection{Synthèse et conclusion}

Les méthodes de Deep Learning surpassent les
approches classiques en précision, en robustesse et en capacité
d'interprétation. Ce sont aujourd'hui les architectures hybrides
CNN/Transformer, ainsi que les modèles multi-vues, qui définissent l'état de
l'art de la classification de la densité mammaire, en combinant analyse de
texture locale, compréhension contextuelle globale et explicabilité
clinique. C'est cette littérature qui a guidé la construction de notre
propre pipeline expérimental, structuré autour des grandes familles
d'architectures suivantes :
\begin{itemize}
    \item extraction de caractéristiques par Vision Transformer,
    \item fine-tuning direct de backbones profonds,
    \item classification hiérarchique,
    \item fusion d'attributs texturaux -- histogramme et GLCM,
    \item comparaison de backbones CNN légers (RexNet150 face à ResNet50 et
    DenseNet121),
    \item intégration d'une approche multi-vue CC/MLO.
\end{itemize}

\section{Correspondance entre nos approches et la littérature}
\label{sec:correspondance}

Chacune des architectures explorées dans ce mémoire s'appuie directement sur
un ou plusieurs travaux cités :

\begin{itemize}
    \item \textbf{Approche 2 (ViT + ResNet50)} reprend le principe hybride
    convolution/attention de MammoViT \cite{ref19}
    (Figure~\ref{fig:mammovit_architecture}), qui combine un extracteur
    ResNet-50 pour les motifs locaux et un Vision Transformer pour le
    contexte global du tissu.
    \item \textbf{Approches 3, 4 et 5 (fusion image + histogramme / GLCM)}
    prolongent la logique des descripteurs texturaux classiques (GLCM
    \cite{ref22}, histogrammes \cite{ref10}) en les combinant à une branche
    convolutive profonde plutôt que de les utiliser seuls.
    \item \textbf{Approche 6 (architectures à double branche)} s'inspire à
    la fois de la comparaison multi-backbone de Pereira et al. \cite{ref9}
    (ResNet50/DenseNet, voir nos variantes \emph{fusion même-vue}) et
    de l'argument clinique de complémentarité CC/MLO détaillé dans notre
    état de l'art (Figures~\ref{fig:vue_mlo_journal} et
    \ref{fig:vue_cc_journal}) : certaines lésions ne sont visibles que sur
    une seule des deux incidences, ce qui justifie la fusion plutôt que le
    choix d'une vue unique.
    \item \textbf{L'idée de classification hiérarchique / en cascade}
    reprend directement le principe de
    Trace4BDensity \cite{trace4bdensity} (trois classifieurs en cascade
    imitant le raisonnement par élimination du radiologue).
    \item \textbf{L'apprentissage fédéré} (Section~\ref{sec:federe})
    répond directement à la contrainte de souveraineté et de
    confidentialité des données médicales identifiée en
    Section~\ref{sec:etat-art-federe} : la plupart des travaux cités
    ci-dessus s'appuient sur une cohorte déjà centralisée, alors qu'en
    pratique les données médicales réelles restent réparties entre
    institutions et ne peuvent pas être regroupées librement, ce qui
    motive une approche par apprentissage fédéré plutôt que centralisé.
\end{itemize}


\begin{figure}[H]
    \centering
    \includegraphics[width=0.65\textwidth]{image/médiolatérale.png}

    \caption{Vue oblique médio-latérale (MLO), reprise de l'état de l'art
    pour rappel du contexte anatomique justifiant l'Approche 6.}
    \label{fig:vue_mlo_journal}
\end{figure}

\begin{figure}[H]
    \centering
     \includegraphics[width=0.65\textwidth]{image/pectoralis.png}
    \caption{Complémentarité des vues CC et MLO, reprise de l'état de l'art.}
    \label{fig:vue_cc_journal}
\end{figure}



% ============================================================================
\chapter{Résultats expérimentaux par approche}
\label{ch:approches}

\section{Cadre méthodologique}

\subsection{Contraintes et ressources}

Le projet s'est déroulé dans un cadre de recherche appliquée soumis à
plusieurs contraintes :
\begin{itemize}
    \item \textbf{Techniques} : manipulation d'images médicales volumineuses
    au format DICOM, anonymisation stricte, et usage exclusif d'outils open
    source ;
    \item \textbf{Matérielles} : serveur GPU (4 cartes de 3\,Go) de la
    Faculté Polytechnique de Mons ;
    \item \textbf{Encadrement} : supervision de M. Xavier Lessage, maître de
    mémoire, avec des échanges réguliers avec le professeur Saïd Mahmoudi.
\end{itemize}

\subsection{Méthodologie et outils}

Le travail a suivi une approche itérative et expérimentale, alternant
conception, entraînement et validation de modèles. Les outils principaux
utilisés étaient :
\begin{itemize}
    \item \textbf{Langages et frameworks} : Python, PyTorch, timm,
    Transformers (Hugging Face), OpenCV, Scikit-image, Scikit-learn,
    Pydicom, Matplotlib/Seaborn ;
    \item \textbf{Suivi et documentation} : GitHub et Overleaf ;
    \item \textbf{Réduction de dimension} : extraction de caractéristiques
    par backbone (vecteurs 768D pour ViT, 2048D pour ResNet50), descripteurs
    statistiques compacts (histogramme 256D, GLCM à 5 valeurs), et t-SNE
    pour la projection de l'espace latent en 2D ;
    \item \textbf{Évaluation} : accuracy, F1-score, matrice de confusion, et
    analyse comparative.
\end{itemize}

\section{Prétraitement et augmentation des données}
\label{sec:augmentation}

VinDr-Mammo présente un déséquilibre de classe très marqué. La classe
Density C (tissu dense hétérogène) représente à elle seule plus des trois
quarts des images, alors que Density A (tissu presque entièrement adipeux)
en représente moins de 0,5\,\%. Un entraînement direct sur cette
distribution conduirait le modèle à ignorer purement et simplement les
classes minoritaires -- c'est d'ailleurs ce qu'on observe dans plusieurs de
nos résultats (voir par exemple Section~\ref{sec:test-seuil-a}). Une
stratégie d'augmentation ciblée sur les classes minoritaires a donc été
mise en place dans l'ensemble des scripts d'entraînement du projet.

\subsection{Vérification préalable : cohérence gauche/droite}

Avant d'appliquer une augmentation quelconque, il fallait s'assurer qu'elle
ne risquait pas d'introduire une incohérence d'étiquetage. Une vérification
a donc été menée sur l'ensemble des études de VinDr-Mammo : pour une même
patiente (même \texttt{study\_id}), la densité mammaire annotée est-elle la
même sur toutes les images, les deux vues CC et MLO comprises ? Le résultat
confirme une cohérence quasi totale : \textbf{4968 études sur 5000
(99,4\,\%)} présentent une densité identique sur l'ensemble de leurs
images. Cette cohérence permet d'appliquer des transformations géométriques
et photométriques sur les images sans risque de mélanger des étiquettes
contradictoires.

\subsection{Transformations appliquées}

Les transformations appliquées visent à augmenter la diversité tout en
respectant l'anatomie mammaire :
\begin{itemize}
    \item rotation légère ($\pm15^{\circ}$) pour simuler les variations
    d'orientation lors de l'acquisition ;
    \item retournement horizontal ($p=0{,}5$) pour augmenter la diversité
    inter-patientes ;
    \item ajustement de luminosité et de contraste ($\pm10\,\%$) pour
    reproduire les différences d'exposition entre appareils ;
    \item recadrage aléatoire (échelle $0{,}9$--$1{,}1$) pour renforcer la
    robustesse aux translations et légers décalages de cadrage.
\end{itemize}
Ces transformations ne sont appliquées qu'au jeu d'entraînement. Elles ne
touchent jamais la validation ni le test, pour ne pas fausser l'évaluation.

\subsection{Répartition avant et après suréchantillonnage}

Le suréchantillonnage utilisé dans le code du projet porte les classes
minoritaires jusqu'à 50\,\% de la taille de la classe majoritaire (voir
Section~\ref{sec:approche5}), plutôt qu'un rééquilibrage complet à
l'identique sur les 4 classes. C'est un compromis qui limite la duplication
excessive des classes les plus rares (voir la discussion sur la
mémorisation, Section~\ref{sec:test-seuil-a}). Les tableaux ci-dessous,
calculés directement sur le CSV d'annotations VinDr-Mammo
(\texttt{Codes/compute\_augmentation\_stats.py}), montrent la répartition
réelle avant et après cette étape.

\begin{table}[h]
\centering
\caption{Répartition des images par densité et par split, avant augmentation}
\begin{tabular}{lrr}
\toprule
Classe & Training (16000 images) & Test (4000 images) \\
\midrule
Density A & 80 & 20 \\
Density B & 1528 & 380 \\
Density C & 12232 & 3060 \\
Density D & 2160 & 540 \\
\bottomrule
\end{tabular}
\end{table}

\begin{table}[h]
\centering
\caption{Répartition du jeu d'entraînement après suréchantillonnage ciblé}
\begin{tabular}{lrr}
\toprule
Classe & Avant & Après \\
\midrule
Density A & 80 & 6116 \\
Density B & 1528 & 6116 \\
Density C & 12232 & 12232 (inchangée, classe majoritaire) \\
Density D & 2160 & 6116 \\
\bottomrule
\end{tabular}
\end{table}

Cette stratégie ramène le ratio entre la classe la plus fréquente et la
classe la plus rare de \textbf{1:152,9} (un déséquilibre brut, bien plus
sévère que sur des jeux de données curés) à \textbf{1:2,0} après
suréchantillonnage. Le jeu de test, lui, n'est jamais modifié : il conserve
la distribution naturelle du dataset (Density A = 20 images sur 4000, soit
0,5\,\%). C'est pour ça que les métriques de test sur cette classe restent
statistiquement fragiles même après un entraînement rééquilibré, avec
seulement 10 images de support sur un sous-échantillon de test de 2000
images pour la vue CC (voir Section~\ref{sec:approche5}).

% ============================================================================
\section{Approche 1 -- Architecture univariée fondée sur les Vision Transformers}
\label{sec:approche1}

\subsection{Principe et chaîne de traitement}
Avant d'explorer les architectures hybrides à plusieurs branches, une
première famille de modèles a servi de point de comparaison : les
\emph{Vision Transformers} (ViT). Ils traitent l'image mammographique dans
son ensemble, via un mécanisme d'auto-attention, plutôt que par empilement
de filtres convolutifs locaux comme un CNN classique.

Le pipeline retenu extrait d'abord un vecteur de caractéristiques de
dimension 768 à partir de l'image prétraitée. Ce vecteur alimente un premier
MLP à quatre sorties, qui donne une probabilité par classe BI-RADS. Pour
limiter les erreurs sur les paires de classes visuellement proches
(notamment B/C et C/D), une seconde étape intervient lorsque les deux
classes les plus probables (\emph{top-2}) sont proches en score : un second
MLP, binaire et spécialisé sur cette paire précise, tranche alors entre les
deux.

\begin{figure}[H]
\centering
\begin{tikzpicture}[
    node distance=0.6cm and 0.8cm,
    stepblock/.style={rectangle, draw=blue!70!black, fill=blue!8, thick, rounded corners=5pt,
                      minimum height=1.1cm, minimum width=3.4cm, align=center, font=\small\sffamily, inner sep=6pt},
    featureblock/.style={rectangle, draw=purple!70!black, fill=purple!10, thick, rounded corners=5pt,
                         minimum height=1.1cm, minimum width=3.6cm, align=center, font=\small\sffamily, inner sep=6pt},
    decisionblock/.style={rectangle, draw=orange!80!black, fill=orange!10, thick, rounded corners=5pt,
                          minimum height=1.1cm, minimum width=3.4cm, align=center, font=\small\sffamily, inner sep=6pt},
    outputblock/.style={rectangle, draw=green!60!black, fill=green!15, thick, rounded corners=5pt,
                        minimum height=1.1cm, minimum width=3.6cm, align=center, font=\small\bfseries\sffamily, inner sep=6pt},
    link/.style={-Stealth, line width=1.2pt, color=blue!80!black}
]
    \node (img) [stepblock, fill=blue!10] {Image DICOM \\ \scriptsize (VinDr-Mammo)};
    \node (prep) [stepblock, right=of img, fill=cyan!10] {Prétraitement \\ \scriptsize (normalisation, rognage)};
    \node (trans) [featureblock, right=of prep] {Backbone Transformer \\ \scriptsize (Twins / CvT -- 768D)};
    \node (mlp4) [stepblock, right=of trans, fill=indigo!10] {MLP 4 classes \\ \scriptsize ($P(y|x)$)};

    \node (top2) [decisionblock, below=1.0cm of img] {Sélection top-2 \\ \scriptsize (cas ambigus)};
    \node (mlpbin) [decisionblock, right=of top2, fill=magenta!10] {MLP binaire \\ \scriptsize (arbitrage ciblé)};
    \node (out) [outputblock, right=of mlpbin] {Prédiction finale \\ \scriptsize BI-RADS (A, B, C, D)};

    \draw [link] (img) -- (prep);
    \draw [link] (prep) -- (trans);
    \draw [link] (trans) -- (mlp4);
    \draw [link] (mlp4.south) -- ++(0,-0.4) -| (top2.north);
    \draw [link] (top2) -- (mlpbin);
    \draw [link] (mlpbin) -- (out);
\end{tikzpicture}
\caption{Chaîne de traitement de l'Approche 1 : extraction de caractéristiques
par backbone Transformer, classification directe à 4 classes, puis
arbitrage binaire ciblé en cas d'ambiguïté entre les deux classes les plus
probables.}
\label{fig:flowchart_approche1}
\end{figure}

\subsubsection{Choix du backbone}
Plusieurs backbones Transformer disponibles dans le projet ont d'abord été
comparés qualitativement, avant de lancer les entraînements complets, selon
leur adéquation attendue à la texture fine du tissu mammaire
(Table~\ref{tab:transformer_selection}). Le \textbf{Twins Transformer},
complété par une évaluation de \textbf{CvT-W24}, a été retenu comme le
compromis le plus prometteur, grâce à son alternance entre attention locale
et attention globale.

\begin{table}[h!]
    \centering
    \small
    \begin{tabular}{lp{8.5cm}c}
    \toprule
    \textbf{Backbone} & \textbf{Pertinence pour la mammographie} & \textbf{Retenu} \\
    \midrule
    ViT-Base & Attention globale uniforme ; tend à lisser les micro-textures fines. & Moyen \\
    DeiT & Convient à de petits volumes de données mais résolution spatiale limitée. & Moyen \\
    PVT & Pyramide multi-échelle, adaptée aux variations d'échelle du tissu dense. & Bon \\
    \textbf{Twins / CvT} & \textbf{Combine convolutions locales et attention globale.} & \textbf{Retenu} \\
    \bottomrule
    \end{tabular}
    \caption{Comparaison qualitative des backbones Transformer envisagés pour l'Approche 1.}
    \label{tab:transformer_selection}
\end{table}

\subsection{Résultats sur le jeu de test VinDr-Mammo}
L'évaluation a porté sur les 1808 images du jeu de test indépendant, avec le
backbone Twins Transformer.

\begin{figure}[H]
    \centering
    \includegraphics[width=0.72\textwidth]{image/twins_mlp4_confusion_matrix.png}
    \caption{Matrice de confusion -- classification directe à 4 classes (MLP4, Twins Transformer).}
    \label{fig:twins_mlp4_matrix}
\end{figure}

\begin{figure}[H]
    \centering
    \includegraphics[width=0.72\textwidth]{image/twins_hierarchical_confusion_matrix.png}
    \caption{Matrice de confusion -- pipeline avec arbitrage binaire ciblé (Twins Transformer).}
    \label{fig:twins_hierar_matrix}
\end{figure}

Le classifieur direct à 4 classes atteint une accuracy globale de
\textbf{76.66\,\%}. Ce chiffre cache cependant une forte disparité entre
classes : la classe C, très majoritaire, est bien reconnue (1091/1376, soit
un rappel de 79.29\,\%), alors que la classe A n'est \textbf{jamais}
prédite correctement (rappel de 0\,\%). L'arbitrage binaire ciblé améliore
le rappel de la classe B (149/172, soit 86.63\,\%), mais fait chuter
l'accuracy globale à 60.57\,\%. Contrairement à l'Approche 7
(Section~\ref{sec:approche7}), cette architecture ne répartit pas les
classes en deux groupes fixes : chaque échantillon reçoit son propre
arbitrage, selon les deux classes les plus probables que lui attribue le
MLP4. C'est précisément ce mécanisme qui pénalise la classe D : son rappel
s'effondre à 6.4\,\% (16/250) dans le pipeline complet, alors que le MLP4
seul obtenait déjà 69.6\,\% sur cette même classe -- ce n'est donc pas
l'arbitrage binaire C/D en lui-même qui est en cause (isolé, il retrouve un
rappel D de 75.2\,\%, Figure~\ref{fig:twins_binary_cd}), mais la
combinaison de la sélection top-2 et de cet arbitrage qui, pour la grande
majorité des vrais cas D, aboutit à une prédiction C.

\begin{figure}[H]
    \centering
    \includegraphics[width=0.55\textwidth]{image/twins_binary_2_3_confusion_matrix.png}
    \caption{Classifieur binaire C vs D évalué isolément (hors pipeline hiérarchique), Twins Transformer.}
    \label{fig:twins_binary_cd}
\end{figure}

Les courbes ROC
one-vs-rest du classifieur 4 classes confirment ce constat, avec une AUC de
0.79 pour A contre 0.93, 0.81 et 0.86 pour B, C et D (Figure~\ref{fig:twins_mlp4_roc}) --
la classe A reste
la plus difficile à isoler, malgré une capacité de discrimination correcte
sur les trois autres.

\begin{figure}[H]
    \centering
    \includegraphics[width=0.72\textwidth]{image/twins_mlp4_roc_curves.png}
    \caption{Courbes ROC one-vs-rest du classifieur 4 classes (MLP4, Twins Transformer).}
    \label{fig:twins_mlp4_roc}
\end{figure}

\subsection{Standardisation photométrique des images DICOM}
L'inspection des en-têtes DICOM de VinDr-Mammo a révélé que certaines
images sont enregistrées avec l'attribut \texttt{PhotometricInterpretation}
à \texttt{MONOCHROME1} (fond clair, tissu sombre) plutôt qu'à
\texttt{MONOCHROME2} (fond sombre, tissu clair). Cette inversion de
polarité, si elle n'est pas corrigée, perturbe directement l'apprentissage
des filtres convolutifs et des mécanismes d'attention : la même densité de
tissu ne produit alors pas la même intensité de pixel selon le format
d'origine.

Un module de détection et d'inversion automatique a donc été ajouté au
prétraitement (\texttt{Codes/preprocessing.py}, fonction
\texttt{read\_dicom}), qui ramène systématiquement toute image vers la
convention à fond noir (\texttt{MONOCHROME2}). Pour vérifier que cette
conversion ne modifie pas l'information utile, mais seulement sa polarité,
une analyse à quatre panneaux a été mise en place : image brute, image
corrigée, profil d'intensité 1D sur la ligne médiane, et histogramme des
intensités. La Figure~\ref{fig:monochrome1_analysis} montre le cas d'une
image \texttt{MONOCHROME1} : le profil 1D confirme une inversion exacte et
symétrique du signal. La Figure~\ref{fig:monochrome2_analysis}, pour une
image déjà en \texttt{MONOCHROME2}, montre au contraire une superposition
parfaite entre le signal brut et le signal après prétraitement -- la preuve
que la conversion n'altère rien quand elle n'est pas nécessaire.

\begin{figure}[H]
    \centering
    \includegraphics[width=0.98\textwidth]{image/figure6_monochrome1.png}
    \caption{Image \texttt{MONOCHROME1} avant/après correction : (1) DICOM brut à fond clair ;
    (2) image corrigée à fond noir ; (3) profil d'intensité 1D montrant l'inversion du signal ;
    (4) histogramme des intensités.}
    \label{fig:monochrome1_analysis}
\end{figure}

\begin{figure}[H]
    \centering
    \includegraphics[width=0.98\textwidth]{image/figure7_monochrome2.png}
    \caption{Image \texttt{MONOCHROME2} avant/après passage dans le prétraitement :
    (1) DICOM brut, déjà à fond noir ; (2) image inchangée ; (3) profil 1D montrant la
    superposition exacte des deux signaux ; (4) histogramme des intensités.}
    \label{fig:monochrome2_analysis}
\end{figure}

En complément, une pondération inversement proportionnelle à la fréquence
de chaque classe a été appliquée à la fonction de perte, pour limiter le
biais vers la classe C déjà observé plus haut.

\subsection{Essai avec CvT-W24 et cohérence gauche/droite}
Un second backbone, \textbf{CvT-W24} (\emph{Convolutional Vision
Transformer, Wide-24}), a été testé pour voir si la réintroduction de
convolutions dans les blocs d'attention améliorait la capture des
micro-textures. Après 50 époques d'entraînement (taux d'apprentissage
$10^{-4}$, batch de 16), le modèle dépasse 90\,\% d'accuracy sur
l'entraînement mais stagne sur le test : c'est un cas de surapprentissage
marqué, le modèle mémorise les particularités des images d'entraînement
sans généraliser à de nouvelles patientes.

Avant de changer de stratégie d'architecture, la cohérence des annotations
de densité a été vérifiée entre le sein gauche et le sein droit d'une même
patiente sur l'ensemble de VinDr-Mammo, avec la même méthode que celle de
la Section~\ref{sec:augmentation} : \textbf{99.4\,\%} des études présentent
une densité identique des deux côtés. Cette cohérence quasi totale confirme
la fiabilité de l'étiquetage du dataset, et justifie l'usage de
transformations symétriques (flip horizontal) dans la stratégie
d'augmentation.

Trois pistes d'amélioration ont été retenues à l'issue de cette première
approche :
\begin{enumerate}
    \item renforcer la régularisation (dropout plus élevé, label smoothing) pour limiter le surapprentissage observé avec CvT-W24 ;
    \item cibler l'augmentation de données sur les classes minoritaires (A et D) plutôt que sur l'ensemble du jeu de données ;
    \item remplacer l'arbitrage binaire local par une véritable architecture hiérarchique en cascade, plus proche du raisonnement clinique par élimination progressive.
\end{enumerate}

\subsection{Bilan et transition}
La comparaison entre classification directe
(Figure~\ref{fig:twins_mlp4_matrix}) et arbitrage binaire ciblé
(Figure~\ref{fig:twins_hierar_matrix}) montre qu'aucune des deux stratégies,
prise isolément, ne traite correctement le déséquilibre naturel des classes
de densité. La première ignore complètement la classe A ; la seconde
corrige partiellement ce biais, mais au prix d'une accuracy globale
nettement plus faible. Ce constat motive le passage à une architecture de
classification hiérarchique en cascade, structurée selon la logique
diagnostique des radiologues : opposer d'abord tissu peu dense et tissu
dense, puis raffiner à l'intérieur de chaque groupe. C'est ce que développe
la suite de ce chapitre.

% ============================================================================
\section{Approche 2 -- Fine-tuning de backbones (ViT et ResNet50), classification directe}
\label{sec:approche2}

Cette approche sert de référence simple : plutôt que de combiner plusieurs
sources d'information, on entraîne directement un backbone profond de bout
en bout sur l'image entière, sans branche auxiliaire. L'idée est d'avoir un
point de comparaison, pour juger si les architectures hybrides introduites
plus loin apportent vraiment un gain, ou si un simple fine-tuning suffit
déjà à capturer l'essentiel du signal.

\subsection{Architecture}
Le backbone (ViT-Base ou ResNet50, tous deux pré-entraînés sur ImageNet) reçoit
l'image mammographique prétraitée en entrée et produit un vecteur de
caractéristiques global $f_\theta(x)$ après pooling. Sa tête de classification
d'origine est remplacée par une couche linéaire projetant directement vers les
quatre classes de densité :
\[
  \hat{y} = \operatorname{softmax}\big(W_c\, f_\theta(x) + b_c\big), \qquad
  W_c \in \mathbb{R}^{4 \times d}, \; b_c \in \mathbb{R}^4
\]
où $d$ est la dimension du vecteur de sortie du backbone ($d{=}768$ pour
ViT-Base, $d{=}2048$ pour ResNet50 avant projection). L'entraînement optimise
une entropie croisée standard :
\[
  \mathcal{L} = -\frac{1}{N}\sum_{i=1}^{N} \log \hat{p}_{i, y_i}
\]

\subsection{ViT}
\textbf{Terminé.} 10 époques, batch size 10, GTX 1080 Ti (environ 2h18
d'entraînement), pour une Train Acc finale de 74.26\,\% -- ce mode
d'entraînement n'a pas de split de validation, voir Section~\ref{sec:approche2}
plus haut. L'évaluation sur le jeu de test indépendant (4000 images, split
\texttt{test} de VinDr-Mammo, jamais vu pendant l'entraînement) donne une
\textbf{Test Acc de 70.00\,\%}.

\begin{table}[h]
\centering
\caption{ViT (Approche 2) -- Rapport de classification sur le test (4000 images)}
\begin{tabular}{lrrrr}
\toprule
Classe & Précision & Rappel & F1-score & Support \\
\midrule
Density A & 0.32 & 0.65 & 0.43 & 20 \\
Density B & 0.32 & 0.90 & 0.48 & 380 \\
Density C & 0.93 & 0.67 & 0.78 & 3060 \\
Density D & 0.56 & 0.73 & 0.64 & 540 \\
\midrule
\textbf{Accuracy globale} & \multicolumn{3}{c}{} & \textbf{0.70} (4000) \\
Moyenne macro & 0.53 & 0.74 & 0.58 & \\
Moyenne pondérée & 0.82 & 0.70 & 0.73 & \\
\bottomrule
\end{tabular}
\end{table}

\begin{figure}[H]
    \centering
    \includegraphics[width=0.8\textwidth]{image/test_Vit_direct.png}
    \caption{ViT (Approche 2) -- résultat du test final.}
    \label{fig:vit_test}
\end{figure}

Ce résultat mérite d'être souligné : c'est la \textbf{seule architecture de
tout ce mémoire à obtenir un rappel significatif sur la classe A} (0.65,
soit 13 images sur 20 correctement reconnues), là où toutes les
architectures à double branche échouent presque totalement sur cette classe
(rappel nul ou quasi nul, voir notamment Section~\ref{sec:test-seuil-a}).
L'explication la plus probable tient à la nature du mécanisme d'attention
du Transformer, qui regarde l'image dans son ensemble plutôt que de
s'appuyer sur des filtres convolutifs locaux optimisés pour détecter des
motifs de texture fine. Un sein presque entièrement adipeux (classe A) se
définit avant tout par l'\emph{absence} de texture dense sur tout le
cliché -- un signal global que l'attention du ViT semble mieux capter que
les branches convolutives des architectures hybrides. La contrepartie est
un rappel très élevé mais peu précis sur la classe B (0.90 de rappel,
seulement 0.32 de précision) : le modèle confond plutôt A et B entre elles
(20 images de B prédites A) qu'avec les classes plus denses, ce qui reste
cohérent avec la proximité visuelle réelle de ces deux catégories.
\subsection{ResNet50}
Un bug a été identifié et corrigé avant l'entraînement final :
\texttt{FineTunedFeatureExtractor} appelait \texttt{forward\_features()}, qui
renvoie la carte de caractéristiques spatiale avant pooling
($\mathbb{R}^{C\times H\times W}$), au lieu du modèle complet qui renvoie les
logits de classification ($\mathbb{R}^4$). Résultat :
\texttt{CrossEntropyLoss} plantait, puisqu'elle attend un vecteur de scores
par classe et non une carte spatiale.

\textbf{Terminé.} Train Acc finale 76.39\,\% (meilleure époque, même remarque
que pour ViT : pas de split de validation pour cette approche, voir plus
haut). L'évaluation sur le jeu de test indépendant (4000 images) donne une
\textbf{Test Acc de 74.00\,\%}.

\begin{table}[h]
\centering
\caption{ResNet50 (Approche 2) -- Rapport de classification sur le test (4000 images)}
\begin{tabular}{lrrrr}
\toprule
Classe & Précision & Rappel & F1-score & Support \\
\midrule
Density A & 0.30 & 0.95 & 0.46 & 20 \\
Density B & 0.48 & 0.84 & 0.61 & 380 \\
Density C & 0.97 & 0.70 & 0.81 & 3060 \\
Density D & 0.46 & 0.90 & 0.61 & 540 \\
\midrule
\textbf{Accuracy globale} & \multicolumn{3}{c}{} & \textbf{0.74} (4000) \\
Moyenne macro & 0.55 & 0.85 & 0.62 & \\
Moyenne pondérée & 0.85 & 0.74 & 0.76 & \\
\bottomrule
\end{tabular}
\end{table}

\begin{figure}[H]
    \centering
    \includegraphics[width=0.8\textwidth]{image/test_resnet50_direct.png}
    \caption{ResNet50 (Approche 2) -- résultat du test final.}
    \label{fig:resnet50_test}
\end{figure}

C'est le résultat le plus marquant de tout ce mémoire sur la classe A :
\textbf{rappel de 0.95} (19 images sur 20 correctement reconnues), largement
au-dessus de tout ce qui a été observé ailleurs -- y compris le ViT de cette
même Approche 2 (0.65). La précision associée reste faible (0.30, 44 fausses
alarmes sur les 63 images prédites A), mais le signal est net : sur un
problème où toutes les architectures à double branche échouent totalement
sur cette classe (Section~\ref{sec:approche6}), un simple fine-tuning direct
de ResNet50 sur l'image entière suffit à la détecter dans la quasi-totalité
des cas. L'explication la plus probable est la même que pour le ViT : sans
segmentation ni branche auxiliaire, le réseau apprend directement à repérer
l'absence de texture dense sur l'ensemble du cliché, un signal global que les
architectures hybrides, focalisées sur la fusion de features locales,
semblent perdre en route. La contrepartie est un rappel de la classe C en
retrait (0.70) : le modèle « sur-route » vers B et D au profit de la
sensibilité sur A, un compromis qui explique l'écart entre l'accuracy globale
modeste (74\,\%) et le macro F1 relativement élevé (0.62) par rapport aux
architectures qui écrasent la classe A pour maximiser l'accuracy sur la
classe C majoritaire.

% ============================================================================
\section{Approche 3 -- ResNet50 + Histogramme (vues MLO)}
\label{sec:approche3}

\subsection{Architecture}
Cette approche introduit le principe de fusion à deux branches qui reviendra
tout au long du mémoire. Une branche image (ResNet50, 1 canal, sortie
projetée à 512 dimensions) extrait les caractéristiques visuelles globales.
Une seconde branche, purement statistique, résume la distribution des
niveaux de gris de l'image sous forme d'un histogramme à 256 classes
$h \in \mathbb{R}^{256}$, projeté par un petit MLP (\texttt{HistogramMLP})
vers un vecteur de 64 dimensions. L'idée est que la densité mammaire dépend
en partie de la proportion de tissu dense (clair) par rapport au tissu
graisseux (sombre) -- une information que l'histogramme capture directement,
sans passer par des convolutions :
\[
  f_{\text{img}}(x) \in \mathbb{R}^{512}, \qquad
  f_{\text{hist}}(h) \in \mathbb{R}^{64}
\]
Les deux vecteurs sont concaténés puis passés à un MLP final de classification :
\[
  z = \big[f_{\text{img}}(x) \,\Vert\, f_{\text{hist}}(h)\big] \in \mathbb{R}^{576},
  \qquad
  \hat{y} = \operatorname{softmax}\big(\text{MLP}_c(z)\big)
\]

\textbf{Terminé.} Validation Acc \textbf{83.36\%}, Test Acc \textbf{85\%}
(2000 images MLO). Point faible : recall nul (0\%) sur Density A (classe très
minoritaire).

\begin{center}
\includegraphics[width=0.8\textwidth]{image/test_hybrid_resNet50_MLO.png}\\
\includegraphics[width=0.8\textwidth]{image/entrainement_hybrid_resNet50_MLO.png}
\end{center}

% ============================================================================
\section{Approche 4 -- RexNet150 + Histogramme (vues CC)}
\label{sec:approche4}
Architecture strictement identique à l'Approche 3 (Section~\ref{sec:approche3}),
seul le backbone image change (RexNet150 au lieu de ResNet50) ainsi que la vue
utilisée (craniocaudale au lieu de médio-latérale oblique). Ce choix permet une
comparaison directe de l'effet du backbone et de la vue à architecture de
fusion constante.

\textbf{Terminé.} Validation Acc \textbf{81.56\%}, Test Acc \textbf{80\%}
(2000 images CC).
\begin{center}
\includegraphics[width=0.8\textwidth]{image/test_Hybride_RexNet150_CC.png}\\
\includegraphics[width=0.8\textwidth]{image/entrainement_Hybride_RexNet150_CC.png}
\end{center}
% ============================================================================
\section{Approche 5 -- Intégration d'une branche GLCM (comparaison de 3 modèles)}
\label{sec:approche5}

\subsection{Motivation et architecture}
L'histogramme des Approches 3 et 4 ne décrit que la distribution globale des
intensités, sans tenir compte de leur agencement spatial. Or la densité
mammaire se traduit aussi par des propriétés de \emph{texture} du tissu,
comme son homogénéité ou sa granularité, qui échappent à un simple
histogramme. La branche GLCM (\emph{Gray Level Co-occurrence Matrix})
remplace donc la branche histogramme par une branche de texture.

La matrice de co-occurrence $P_{d,\theta}(i,j)$, introduite par Haralick et
al.\ \cite{ref_haralick}, mesure, pour une distance $d$ et une orientation
$\theta$ données, la fréquence à laquelle un pixel d'intensité $i$ se
retrouve voisin d'un pixel d'intensité $j$. Cinq des descripteurs
statistiques proposés dans cet article sont calculés à partir de cette
matrice, formant le vecteur de texture utilisé en entrée du
\texttt{GLCMDescriptorMLP} :
\begin{align}
  \text{Contraste} &= \sum_{i,j} (i-j)^2\, P(i,j) \\
  \text{Homogénéité} &= \sum_{i,j} \frac{P(i,j)}{1+(i-j)^2} \\
  \text{Énergie} &= \sum_{i,j} P(i,j)^2 \\
  \text{Corrélation} &= \sum_{i,j} \frac{(i-\mu_i)(j-\mu_j)\,P(i,j)}{\sigma_i \sigma_j} \\
  \text{Entropie} &= -\sum_{i,j} P(i,j)\, \log P(i,j)
\end{align}
Ces cinq valeurs sont projetées par un MLP léger vers un vecteur de 16
dimensions $f_{\text{glcm}} \in \mathbb{R}^{16}$, concaténé à la branche image
($\mathbb{R}^{512}$) exactement comme pour l'histogramme :
\[
  z = \big[f_{\text{img}}(x) \,\Vert\, f_{\text{glcm}}\big] \in \mathbb{R}^{528}
\]

\subsection{Déséquilibre de classes et pondération}
Comme Density A représente moins de 1\% des données, la loss d'entropie
croisée est pondérée par l'inverse de la fréquence de chaque classe -- une
pondération courante quand le déséquilibre est marqué :
\[
  w_c = \frac{N}{K \cdot n_c}, \qquad
  \mathcal{L} = -\frac{1}{N}\sum_{i=1}^N w_{y_i} \log \hat{p}_{i,y_i}
\]
où $N$ est le nombre total d'exemples, $K{=}4$ le nombre de classes et $n_c$
le nombre d'exemples de la classe $c$. La classe minoritaire est aussi
ré-échantillonnée dans le jeu d'entraînement jusqu'à atteindre 50\% de la
taille de la classe majoritaire, un compromis qui évite à la fois la rareté
extrême de la classe A et une duplication excessive qui favoriserait le
surapprentissage. Le taux d'apprentissage suit une décroissance en cosinus :
\[
  \eta_t = \eta_{\min} + \tfrac{1}{2}(\eta_{\max}-\eta_{\min})
  \Big(1+\cos\big(\tfrac{t}{T_{\max}}\pi\big)\Big)
\]

Un bug de normalisation a aussi été trouvé dans
\texttt{GLCMDescriptorMLP.forward()} : le tenseur d'entrée, déjà normalisé
(valeurs approximativement dans $[-2, 2]$ après la normalisation ImageNet),
était traité comme des pixels bruts $[0, 255]$ avant le calcul du GLCM.
Cela provoquait un débordement \texttt{uint8}, et donc un calcul de texture
sur des valeurs corrompues plutôt que sur l'image réelle. Corrigé dans les
deux copies du projet concernées (\texttt{hybrid\_train\_cc\_rexnet/},
\texttt{hybrid~train~mlo~resnet/}).

\subsection{Modèle 1 -- RexNet150 + GLCM (vues CC)}
\textbf{Terminé (10 époques).} Meilleur modèle retenu à l'Époque 10/10 :
Train Acc 87.98\%, \textbf{Val Acc 73.42\%}. \textbf{Test Acc 74.85\%}
(2000 images CC) -- cohérence val/test qui confirme l'absence de fuite de
données.

\begin{table}[h]
\centering
\caption{RexNet150+GLCM/CC -- Rapport de classification sur le test (2000 images)}
\begin{tabular}{lrrrr}
\toprule
Classe & Précision & Rappel & F1-score & Support \\
\midrule
Density A & 0.33 & 0.10 & 0.15 & 10 \\
Density B & 0.51 & 0.87 & 0.65 & 190 \\
Density C & 0.95 & 0.71 & 0.82 & 1530 \\
Density D & 0.45 & 0.88 & 0.60 & 270 \\
\midrule
\textbf{Accuracy globale} & \multicolumn{3}{c}{} & \textbf{0.75} (2000) \\
Moyenne macro & 0.56 & 0.64 & 0.55 & \\
Moyenne pondérée & 0.84 & 0.75 & 0.77 & \\
\bottomrule
\end{tabular}
\end{table}

\begin{center}
\includegraphics[width=0.6\textwidth]{image/courbe_accuracy_modele1_rexnet150_glcm_cc.png}\\
\includegraphics[width=0.6\textwidth]{image/entrainement_modele1_rexnet150_glcm_cc.png}\\
\includegraphics[width=0.6\textwidth]{image/test_modele1_rexnet150_glcm_cc.png}
\end{center}
L'évaluation dédiée du split de validation donne 73.64\% (1597 images),
un chiffre cohérent avec le 73.42\% rapporté pendant l'entraînement -- la
légère différence vient du \texttt{drop\_last} sur le \texttt{DataLoader}
d'entraînement. Pas de courbe de loss possible pour ce run : elle n'a pas
été enregistrée à l'époque, mais c'est corrigé dans le code pour les
prochains entraînements (voir note méthodologique en fin de section).

\subsubsection{Test du seuil de confiance sur la classe A}
\label{sec:test-seuil-a}

Lors d'un point d'avancement, mon  encadrant a proposé une piste pour gérer la classe A, très minoritaire : si le modèle n'est confiant sur
aucune des classes B, C ou D (par exemple si leur probabilité combinée
reste sous un seuil de 40\,\%), on pourrait classer automatiquement l'image
en A par élimination, plutôt que de laisser le modèle trancher directement
entre les 4 classes. vérifier cette hypothèse en envoyant au modèle des images dont la vraie classe est A, et en regardant les
probabilités de sortie plutôt que la seule prédiction finale.

Ce test a été fait sur les 10 images de classe A du jeu de test (vues CC),
avec le modèle RexNet150+GLCM. Le résultat ne confirme pas l'hypothèse
de départ :

\begin{table}[h]
\centering
\caption{Comportement du modèle sur les 10 images réelles de classe A (test)}
\begin{tabular}{lr}
\toprule
Classe A prédite correctement & 1/10 \\
Confiant ($\geq 40\,\%$) mais faux sur B, C ou D & 9/10 \\
Incertain (aucune classe B/C/D $\geq 40\,\%$) & 0/10 \\
\bottomrule
\end{tabular}
\end{table}

Sur les 9 images mal classées, le modèle attribue presque toujours plus de
96\,\% de probabilité à la classe B, parfois jusqu'à 100,0\,\%. Il n'est
donc jamais réellement \og{}incertain\fg{} face à une image de classe A :
au contraire, il est très confiant, mais dans la mauvaise classe. Cette
observation invalide l'heuristique du seuil de confiance telle que proposée
au départ, qui supposait une zone d'hésitation exploitable entre les
classes. Malgré la pondération de la loss et le suréchantillonnage déjà en
place pour cette classe, le modèle semble avoir surtout appris à ignorer la
classe A au profit de B. La raison la plus probable : les 10 exemples de
test, et le nombre restreint d'exemples natifs disponibles à
l'entraînement avant suréchantillonnage, ne suffisent pas à lui faire
apprendre une frontière de décision fiable. Ce constat pointe donc moins
vers une correction simple par seuillage que vers le besoin d'une vraie
collecte ou augmentation de données pour cette classe.

Ce n'est pas un bug dans le code d'entraînement, mais une limite
structurelle : le suréchantillonnage actuel duplique un très petit nombre
d'images natives (le \texttt{target\_count} fixé à 50\,\% de la classe
majoritaire implique un facteur de duplication de l'ordre de 50 à 75 pour
la classe A) plutôt que d'apporter une réelle diversité. Cela favorise la
mémorisation de ces quelques images plutôt que l'apprentissage d'une notion
générale de tissu peu dense. Deux pistes d'amélioration restent
envisageables, mais n'ont pas pu être mises en \oe{}uvre faute de temps
avant l'échéance :\\
\begin{itemize}
    \item remplacer la pondération de la loss par une \emph{focal loss},
    qui force explicitement le modèle à se concentrer sur les exemples
    difficiles plutôt que sur les exemples majoritaires déjà bien appris ;
    \item appliquer un ré-échantillonnage de type SMOTE directement dans
    l'espace des caractéristiques (et non sur les pixels), comme le fait
    MammoViT \cite{ref19} déjà cité dans notre état de l'art, plutôt qu'une
    simple duplication d'images.
\end{itemize}
\todo{reproduire ce même test du seuil de confiance sur un autre modèle une
fois disponible, pour déterminer si cet échec sur la classe A est propre à
cette architecture (RexNet150+GLCM/CC) ou un phénomène général du projet}

\subsection{Modèle 2 -- ResNet50 + GLCM (vues MLO)}
\textbf{Terminé (10 époques).} Meilleur modèle retenu à l'Époque 10/10 :
Train Acc 85.78\%, \textbf{Val Acc 72.81\%}. \textbf{Test Acc 75.95\%}
(2000 images MLO) -- de nouveau, l'écart val/test n'est que de 3 points, ce
qui confirme l'absence de fuite de données pour ce second modèle aussi.
L'évaluation dédiée du split de validation confirme 72.86\,\% (1603
images), un résultat cohérent avec le 72.81\,\% rapporté pendant
l'entraînement.

\begin{figure}[H]
    \centering
    \includegraphics[width=0.8\textwidth]{image/entrainement_modele2_resnet50_glcm_mlo.png}
    \caption{ResNet50+GLCM (vues MLO) -- résultat de l'entraînement (matrice de confusion sur validation, meilleure époque).}
    \label{fig:glcm_mlo_entrainement}
\end{figure}

La matrice de validation montre un profil très proche de celle du test.
La classe A reste difficile (3 succès sur 13, rappel 0.23), mais fait
légèrement mieux qu'en test, où seul 1 succès sur 10 était obtenu -- un
écart plausible vu le petit nombre d'échantillons des deux côtés. Le reste
de la confusion suit le même schéma que partout ailleurs dans ce mémoire :
les erreurs se concentrent presque uniquement entre classes voisines (B
confondue avec C, C avec D), jamais entre classes opposées comme A et D,
ce qui reste cohérent avec la nature progressive de la densité mammaire.

\begin{table}[h]
\centering
\caption{ResNet50+GLCM/MLO -- Rapport de classification sur le test (2000 images)}
\begin{tabular}{lrrrr}
\toprule
Classe & Précision & Rappel & F1-score & Support \\
\midrule
Density A & 0.33 & 0.10 & 0.15 & 10 \\
Density B & 0.49 & 0.92 & 0.64 & 190 \\
Density C & 0.95 & 0.73 & 0.83 & 1530 \\
Density D & 0.49 & 0.84 & 0.62 & 270 \\
\midrule
\textbf{Accuracy globale} & \multicolumn{3}{c}{} & \textbf{0.76} (2000) \\
Moyenne macro & 0.56 & 0.65 & 0.56 & \\
Moyenne pondérée & 0.84 & 0.76 & 0.78 & \\
\bottomrule
\end{tabular}
\end{table}

\begin{figure}[H]
    \centering
    \includegraphics[width=0.8\textwidth]{image/courbe_accuracy_ modele2_resnet50_glcm_mlo.png}
    \caption{ResNet50+GLCM (vues MLO) -- évolution de l'accuracy sur les 10 époques.}
    \label{fig:glcm_mlo_courbe}
\end{figure}

La courbe montre la même dynamique que le Modèle 1 (CC) : le Train Acc
progresse régulièrement, la validation est plus instable, avec un net
décrochage à l'Époque 3 (53.19\,\%, en retrait par rapport aux 57.06\,\% de
l'Époque 2), puis un vrai saut à l'Époque 5 (71.31\,\%) qui marque le début
réel de la convergence. Les 5 premières époques servent surtout au modèle à
sortir d'une zone de réglages instables avant de trouver une direction
stable.

\begin{figure}[H]
    \centering
    \includegraphics[width=0.8\textwidth]{image/test_ modele2_resnet50_glcm_mlo.png}
    \caption{ResNet50+GLCM (vues MLO) -- résultat du test final.}
    \label{fig:glcm_mlo_test}
\end{figure}

Ce second modèle confirme, sur une vue différente (MLO au lieu de CC) et un
backbone différent (ResNet50 au lieu de RexNet150), exactement le même
profil d'erreurs que le Modèle 1. La classe A reste très fragile (1 succès
sur 10, précision 0.33), l'essentiel de la confusion se concentre entre
classes voisines (B/C et C/D), et le résultat global (75.95\,\%) est même
un peu supérieur à celui du Modèle 1 sur CC (74.85\,\%). Cette cohérence
entre deux combinaisons backbone/vue différentes renforce l'idée que la
difficulté sur la classe A est structurelle -- liée à la rareté des données,
voir Section~\ref{sec:test-seuil-a} -- plutôt que propre à une architecture
ou une vue en particulier.

\subsection{Modèle 3 -- RexNet150 + Histogramme (vues MLO)}
\label{sec:modele3-invalide}

\textbf{Historique du bug.} La relecture du code source a révélé que le
premier run réalisé pour ce modèle (Train Acc 84.98\%, Val Acc 86.68\%,
Test Acc 84\%) chargeait ses données d'entraînement et le checkpoint
pré-entraîné de la branche image depuis un répertoire personnel
appartenant à un autre compte du cluster, via des valeurs codées en dur
dans la configuration du script (\texttt{ANNOTATIONS\_CSV},
\texttt{IMAGE\_ROOT}, \texttt{FINETUNED\_WEIGHTS\_PATH}). Ce résultat ne
constituait donc pas une réplication indépendante ; il a été retiré du
mémoire, et le code corrigé pour pointer vers les données VinDr-Mammo
partagées et un checkpoint entraîné localement.

\smallskip
\noindent\textit{Note méthodologique (log/loss) :} un
\texttt{history["train\_loss"]}/\texttt{history["val\_loss"]} a été ajouté aux
scripts d'entraînement (GLCM, fédéré) afin de pouvoir tracer les courbes de
perte train/validation sur les prochains entraînements, en complément des
courbes d'accuracy.

\textbf{Réentraînement avec Focal Loss.} Une fois le code corrigé, ce
modèle a été réentraîné avec les données propres au projet. Comme la
pondération de classe simple s'était déjà montrée insuffisante pour la
classe A ailleurs dans ce mémoire (Section~\ref{sec:test-seuil-a}), ce
réentraînement utilise la \emph{Focal Loss} \cite{ref_focal_loss}
plutôt qu'une simple \texttt{CrossEntropyLoss} pondérée :
\[
  \mathcal{L}_{\text{focal}} = -\sum_i (1 - \hat{p}_{i,y_i})^\gamma \log \hat{p}_{i,y_i}, \qquad \gamma = 2
\]
qui réduit dynamiquement le poids des exemples déjà bien classés (quelle
que soit leur classe), pour concentrer l'apprentissage sur les cas
difficiles -- notamment les rares exemples natifs de la classe A.

\textbf{Terminé (10 époques).} Train Acc 79.25\,\%, \textbf{Val Acc
83.14\,\%} (meilleure époque, retenue).

\begin{figure}[H]
    \centering
    \includegraphics[width=0.8\textwidth]{image/entrainement_modele3_focal.png}
    \caption{RexNet150 + Histogramme (vues MLO), réentraîné avec Focal Loss -- résultat de l'entraînement (matrice de confusion sur validation, meilleure époque).}
    \label{fig:modele3_focal_entrainement}
\end{figure}

\textbf{Test Acc 84.00\,\%} (2000 images MLO) -- l'un des meilleurs
résultats de tout ce mémoire, proche de l'Approche 3 (85\,\%) et de la
variante siamoise de l'Approche 6 (85\,\%).

\begin{figure}[H]
    \centering
    \includegraphics[width=0.8\textwidth]{image/test_modele3_focal.png}
    \caption{RexNet150 + Histogramme (vues MLO), réentraîné avec Focal Loss -- résultat du test final.}
    \label{fig:modele3_focal_test}
\end{figure}

Constat honnête à souligner : malgré la Focal Loss, le \textbf{rappel de
la classe A reste à 0\,\%} (0/10 en test, 0/9 en validation) -- aucune
amélioration sur cette classe précise. Ça confirme ce qui était déjà
suspecté : le problème de la classe A n'est pas une question de
pondération de la loss (simple ou dynamique), mais de rareté brute des
données. Avec seulement une poignée d'exemples natifs disponibles à
l'entraînement, aucune technique de repondération ne peut faire apprendre
au modèle une frontière de décision fiable pour cette classe -- il
faudrait une vraie collecte ou augmentation de données, pas un changement
de fonction de perte.

Le reste de la matrice montre un profil plus contrasté que les autres
modèles du mémoire : la classe C est excellente (rappel 0.96), la classe
B correcte (rappel 0.64), mais la classe D chute à un rappel de
\textbf{0.33} seulement -- la majorité des vraies images D (181 sur 270)
sont classées C. C'est le rappel D le plus faible observé sur une
architecture de ce type dans ce mémoire, ce qui suggère que la Focal Loss,
en poussant le modèle à mieux traiter les cas difficiles en général, a pu
déplacer une partie de la confusion de B/C vers C/D plutôt que de la
réduire uniformément -- un compromis à garder en tête plutôt qu'une
amélioration inconditionnelle.

\subsection{Vérification de la cohérence entre validation et test}
\label{sec:coherence-val-test}

Pour ce modèle, comme pour tous les modèles hybrides de ce mémoire, le
split train/validation est réalisé \emph{avant} toute étape de
suréchantillonnage (voir Section~\ref{sec:augmentation}). C'est un choix
méthodologique important : si le ré-échantillonnage de la classe
minoritaire était appliqué avant la séparation train/validation, les mêmes
images dupliquées se retrouveraient des deux côtés du split, ce qui
gonflerait artificiellement la Val Acc sans refléter une vraie capacité de
généralisation.

Cette précaution se vérifie directement dans nos résultats : la Val Acc
(73.42\%) et la Test Acc (74.85\%) restent proches l'une de l'autre, avec
un écart de seulement 1,4 point. Cette cohérence est le signe que le
modèle généralise correctement à des données réellement non vues, plutôt
que de mémoriser des exemples dupliqués présents des deux côtés du split.

% ============================================================================
\section{Approche 6 -- Architecture à double branche}
\label{sec:approche6}

Quatre architectures sont comparées ici. La première suit une architecture
multi-vue classique : deux vues différentes (CC et MLO), chacune traitée
par un backbone dédié -- un principe courant dans la littérature sur la
classification mammographique multi-vue \cite{ref9, ref14}. Les trois
suivantes explorent d'autres variantes présentes dans cette même
littérature, pour voir si l'architecture originale est optimale ou si
d'autres choix de conception feraient mieux sur les mêmes données.

\subsection{Appariement des vues CC et MLO}

Une architecture à double branche image suppose de disposer, pour un même
sein, à la fois d'une vue CC et d'une vue MLO. Cette contrainte n'est pas
garantie automatiquement par un jeu de données brut. Notre implémentation
(\texttt{HybridMammographyDataset}) regroupe donc systématiquement les
images par \texttt{(study\_id, laterality)}, et ne conserve que les paires
où les deux vues sont effectivement présentes.

Appliqué à l'ensemble des 20000 images de VinDr-Mammo, cet appariement donne
le résultat suivant :
\begin{itemize}
    \item \textbf{Images conservées} (paires CC+MLO complètes) : 19998
    \item \textbf{Images écartées} (vue manquante) : 2
    \item \textbf{Paires study/laterality retenues} : 9999 (9999 vues CC, 9999
    vues MLO)
\end{itemize}
VinDr-Mammo s'avère ici quasiment parfaitement structuré en paires CC/MLO
complètes, conformément à son protocole d'acquisition standard à deux vues
par sein : seules 2 images sur 20000 sont concernées par ce filtrage.
L'entraînement de l'Approche 6 s'appuie donc directement sur la
quasi-totalité du jeu de données disponible, sans perte significative liée
à l'appariement.

\subsection{Architecture originale -- vues hétérogènes (RexNet150/CC + ResNet50/MLO)}
Deux branches image indépendantes, chacune spécialisée sur une vue et un
backbone différents, sont concaténées avant classification :
\[
  z = \big[f^{\text{RexNet150}}_{\text{CC}}(x_{\text{CC}})
  \,\Vert\, f^{\text{ResNet50}}_{\text{MLO}}(x_{\text{MLO}})\big]
  \in \mathbb{R}^{1024}
\]
La stratégie retenue consiste à pré-entraîner chaque branche séparément sur
sa vue respective (Approche 4.5), puis à geler ces poids lors de
l'entraînement du modèle fusionné. Cela stabilise l'optimisation conjointe
et limite le risque de surapprentissage sur un jeu de données de taille
modérée.

Quatre bugs réels ont été trouvés et corrigés dans
\texttt{hybrid~2~branches~images/}. Le gel de branche était silencieusement
inopérant : le code testait l'existence d'attributs au mauvais nom via
\texttt{hasattr}. Le chargement des poids pré-entraînés était conditionné à
tort sur un seul argument \texttt{--backbone} de la ligne de commande, si
bien qu'une seule branche sur deux chargeait effectivement ses poids.
L'initialisation des poids de la couche de classification était placée
après un \texttt{return} dans le code, donc jamais exécutée. Et le jeu de
test n'appariait jamais réellement les images CC et MLO d'un même examen :
il retombait systématiquement sur la même image dupliquée pour les deux
branches, faute d'un mécanisme d'appariement par identifiant d'étude.

\textbf{Terminé (10 époques).} Train Acc 79.74\,\%, Val Acc 83.24\,\% en fin
d'entraînement, mais le meilleur modèle est retenu dès l'\textbf{Époque
2/10} (Val Acc 85.80\,\%) -- c'est ce checkpoint, et non celui de la
dernière époque, qui sert de référence pour la suite.

\begin{figure}[H]
    \centering
    \includegraphics[width=0.8\textwidth]{image/courbe_accuracy_architecture_originale.png}
    \caption{Architecture originale (RexNet150/CC + ResNet50/MLO) -- évolution de l'accuracy sur les 10 époques.}
    \label{fig:archi_originale_courbe}
\end{figure}

La validation atteint son maximum dès l'Époque 2, puis oscille entre
82,95\,\% et 84,66\,\% sans jamais redépasser ce premier pic, alors que
l'entraînement continue lentement à progresser jusqu'à la fin. C'est un
signe de surapprentissage précoce plutôt discret : le modèle n'explose pas
en validation, mais il ne s'améliore plus non plus une fois les branches
pré-entraînées gelées et seule la couche de fusion encore mise à jour.

\begin{figure}[H]
    \centering
    \includegraphics[width=0.8\textwidth]{image/entrainement_architecture_originale.png}
    \caption{Architecture originale (RexNet150/CC + ResNet50/MLO) -- résultat de l'entraînement (matrice de confusion sur validation, meilleure époque).}
    \label{fig:archi_originale_entrainement}
\end{figure}

Sur les 352 images de validation de ce split, aucune n'appartient à la
classe A -- un effet du filtrage par paire CC/MLO combiné à la rareté déjà
connue de cette classe dans VinDr-Mammo. Le reste de la matrice suit un
schéma familier : les erreurs se concentrent entre classes voisines (B/C,
C/D), avec 17 images de classe C classées D et 14 images de classe D
classées C.

\textbf{Terminé.} \textbf{Test Acc 83.00\,\%} (2000 paires CC+MLO du test),
cohérent avec le Val Acc de 85.80\,\% retenu à l'Époque 2 (écart de moins de
3 points, sans signe d'un effondrement en généralisation malgré le
surapprentissage précoce discret déjà relevé sur la courbe de validation).

\begin{table}[h]
\centering
\caption{Architecture originale (RexNet150/CC + ResNet50/MLO) -- Rapport de classification (test, 2000 paires)}
\begin{tabular}{lrrrr}
\toprule
Classe & Précision & Rappel & F1-score & Support \\
\midrule
Density A & 0.14 & 0.20 & 0.17 & 10 \\
Density B & 0.72 & 0.57 & 0.64 & 190 \\
Density C & 0.91 & 0.88 & 0.89 & 1530 \\
Density D & 0.57 & 0.76 & 0.65 & 270 \\
\midrule
\textbf{Accuracy globale} & \multicolumn{3}{c}{} & \textbf{0.83} (2000) \\
Moyenne macro & 0.59 & 0.60 & 0.59 & \\
Moyenne pondérée & 0.84 & 0.83 & 0.83 & \\
\bottomrule
\end{tabular}
\end{table}

\begin{figure}[H]
    \centering
    \includegraphics[width=0.8\textwidth]{image/test_architecture_originale.png}
    \caption{Architecture originale (RexNet150/CC + ResNet50/MLO) -- résultat du test final.}
    \label{fig:archi_originale_test}
\end{figure}

Avec un rappel de 0.20 sur la classe A (2 succès sur 10), cette variante fait
mieux que les architectures à fusion même-vue vues plus haut (0\,\% pour la
paire MLO). Le reste de la matrice reste cohérent avec le schéma habituel :
les erreurs se concentrent sur les paires de classes adjacentes (B/C, C/D),
pas de confusion aberrante entre classes éloignées.

\subsection{Variante siamoise -- poids partagés (EfficientNet-B0)}
Plutôt que deux backbones distincts, cette variante applique un unique
backbone EfficientNet-B0 aux deux vues, avec des poids partagés (\emph{tied
weights}) :
\[
  z = \big[f_\theta(x_{\text{CC}}) \,\Vert\, f_\theta(x_{\text{MLO}})\big],
  \qquad \theta \text{ identique pour les deux passes}
\]
Ce choix de conception, courant dans la littérature multi-vue en
mammographie, repose sur l'hypothèse que les deux vues partagent une même
distribution de texture bas niveau. Il permet aussi de diviser par deux le
nombre de paramètres du modèle par rapport à l'architecture hétérogène.
\textbf{Terminé (10 époques).} Train Acc 86.60\,\%, Val Acc 81.53\,\% en fin
d'entraînement, meilleur modèle retenu à l'\textbf{Époque 4/10} (Val Acc
84.38\,\%).

\begin{figure}[H]
    \centering
    \includegraphics[width=0.8\textwidth]{image/courbe_accuracy_siamoise.png}
    \caption{Variante siamoise (EfficientNet-B0, poids partagés) -- évolution de l'accuracy sur les 10 époques.}
    \label{fig:siamoise_courbe}
\end{figure}

Contrairement à l'architecture hétérogène, ici les deux branches partagent
les mêmes poids et continuent d'apprendre ensemble tout au long de
l'entraînement -- le Train Acc progresse de façon quasi continue jusqu'à
86,60\,\%. La validation, elle, culmine tôt (Époque 4) puis redescend
progressivement, avec un creux marqué à l'Époque 8 (80,11\,\%). L'écart
grandissant entre les deux courbes en fin d'entraînement est un signe de
surapprentissage plus net que pour l'architecture originale, cohérent avec
le fait que ce modèle continue réellement à apprendre (poids partagés non
gelés) plutôt que de se limiter à une couche de fusion.

\begin{figure}[H]
    \centering
    \includegraphics[width=0.8\textwidth]{image/entrainement_siamoise.png}
    \caption{Variante siamoise (EfficientNet-B0, poids partagés) -- résultat de l'entraînement (matrice de confusion sur validation, meilleure époque).}
    \label{fig:siamoise_entrainement}
\end{figure}

Même absence de classe A que pour l'architecture originale sur ce split de
352 images -- attendu, puisque c'est le même jeu de validation. La
confusion se concentre à nouveau sur les paires de classes adjacentes,
avec toutefois un peu moins d'erreurs C/D (4 seulement classées D à tort)
que pour l'architecture hétérogène.

\textbf{Terminé.} \textbf{Test Acc 85.00\,\%} (2000 paires CC+MLO du test) --
le \textbf{meilleur résultat des quatre variantes de l'Approche 6}, et
légèrement au-dessus même du Val Acc retenu (84.38\,\% à l'Époque 4).

\begin{table}[h]
\centering
\caption{Variante siamoise (EfficientNet-B0, poids partagés) -- Rapport de classification (test, 2000 paires)}
\begin{tabular}{lrrrr}
\toprule
Classe & Précision & Rappel & F1-score & Support \\
\midrule
Density A & 0.12 & 0.20 & 0.15 & 10 \\
Density B & 0.69 & 0.59 & 0.64 & 190 \\
Density C & 0.89 & 0.92 & 0.91 & 1530 \\
Density D & 0.69 & 0.61 & 0.65 & 270 \\
\midrule
\textbf{Accuracy globale} & \multicolumn{3}{c}{} & \textbf{0.85} (2000) \\
Moyenne macro & 0.60 & 0.58 & 0.59 & \\
Moyenne pondérée & 0.84 & 0.85 & 0.84 & \\
\bottomrule
\end{tabular}
\end{table}

\begin{figure}[H]
    \centering
    \includegraphics[width=0.8\textwidth]{image/test_siamoise.png}
    \caption{Variante siamoise (EfficientNet-B0, poids partagés) -- résultat du test final.}
    \label{fig:siamoise_test}
\end{figure}

Même rappel sur la classe A que l'architecture originale (0.20, 2 succès sur
10), mais une accuracy globale supérieure (85\,\% contre 83\,\%), portée
par une meilleure reconnaissance de la classe C majoritaire (rappel 0.92
contre 0.88). Le partage de poids entre les deux vues, malgré le
surapprentissage plus net observé en fin d'entraînement (Section
précédente), n'empêche donc pas cette variante de mieux généraliser que
l'architecture à branches hétérogènes sur le jeu de test final.

\subsection{Fusion même-vue -- RexNet150 + ResNet50 (vues MLO)}
Contrairement aux variantes précédentes qui combinent deux \emph{vues}
différentes, cette architecture applique deux backbones différents à la
\emph{même} image :
\[
  z = \big[f^{\text{RexNet150}}(x_{\text{MLO}})
  \,\Vert\, f^{\text{ResNet50}}(x_{\text{MLO}})\big]
\]
Ce schéma, connu dans la littérature sous le nom de \emph{Fusion of Hybrid
Deep Features} (FHDF), part du principe que des architectures différentes
capturent des biais inductifs complémentaires sur une même entrée -- par
exemple, les connexions résiduelles de ResNet50 favorisent des motifs
différents des convolutions rectifiées de RexNet150. Leur combinaison peut
donc être bénéfique même sans information de vue supplémentaire.

\textbf{Terminé (10 époques).} Train Acc 88.74\,\%, Val Acc 74.62\,\%.
\textbf{Test Acc 74.85\,\%} (2000 images MLO) : la cohérence entre
validation et test est raisonnable, sans signe d'un écart anormal.

\begin{figure}[H]
    \centering
    \includegraphics[width=0.8\textwidth]{image/courbe_accuracy_fusion_mlo_rexnet150_resnet50.png}
    \caption{RexNet150+ResNet50 (fusion même-vue, MLO) -- évolution de l'accuracy sur les 10 époques.}
    \label{fig:fusion_mlo_courbe}
\end{figure}

La courbe d'entraînement mérite un commentaire à part. La progression de
l'accuracy sur le jeu d'entraînement est parfaitement régulière, ce qu'on
attend d'un modèle qui apprend correctement. La validation, elle, zigzague
fortement d'une époque à l'autre : elle retombe par exemple à 55\,\% à
l'Époque 2, remonte à 71\,\% à l'Époque 3, puis rechute à 55\,\% à l'Époque
4. Ce comportement n'est pas synonyme d'un entraînement instable ou raté.
Avec seulement 1600 images de validation, et un \textit{scheduler} en
cosinus qui fait encore varier le taux d'apprentissage de façon marquée sur
les premières époques, il est courant que le modèle traverse plusieurs
configurations de poids temporairement moins bonnes avant de converger. Ce
qui compte, c'est la tendance de fond, nettement à la hausse, et le fait
que le dernier point (Époque 10, 74.62\,\%) soit aussi le meilleur de toute
la série. Le modèle n'était donc probablement pas encore arrivé en fin de
convergence quand l'entraînement s'est arrêté, et un nombre d'époques plus
élevé aurait sans doute continué à améliorer un peu le résultat.

\begin{figure}[H]
    \centering
    \includegraphics[width=0.8\textwidth]{image/entrainement_fusion_mlo_rexnet150_resnet50.png}
    \caption{RexNet150+ResNet50 (fusion même-vue, MLO) -- résultat de l'entraînement (matrice de confusion sur validation, meilleure époque).}
    \label{fig:fusion_mlo_entrainement}
\end{figure}

La matrice de confusion de validation confirme ce qu'on observait déjà sur
le test. La classe A reste presque invisible pour le modèle (2 succès sur
13 échantillons), et la confusion se concentre surtout entre les classes B
et C d'un côté, C et D de l'autre. Ce sont des paires de classes
visuellement proches, ce qui reste cohérent avec la nature continue de la
densité mammaire plutôt que quatre catégories parfaitement séparables.

\begin{table}[h]
\centering
\caption{RexNet150+ResNet50 (fusion même-vue, MLO) -- Rapport de classification (test, 2000 images)}
\begin{tabular}{lrrrr}
\toprule
Classe & Précision & Rappel & F1-score & Support \\
\midrule
Density A & 0.00 & 0.00 & 0.00 & 10 \\
Density B & 0.50 & 0.87 & 0.64 & 190 \\
Density C & 0.95 & 0.72 & 0.82 & 1530 \\
Density D & 0.46 & 0.87 & 0.60 & 270 \\
\midrule
\textbf{Accuracy globale} & \multicolumn{3}{c}{} & \textbf{0.75} (2000) \\
Moyenne macro & 0.48 & 0.61 & 0.51 & \\
Moyenne pondérée & 0.84 & 0.75 & 0.77 & \\
\bottomrule
\end{tabular}
\end{table}

\begin{figure}[H]
    \centering
    \includegraphics[width=0.8\textwidth]{image/test_fusion_mlo_rexnet150_resnet50.png}
    \caption{RexNet150+ResNet50 (fusion même-vue, MLO) -- résultat du test final.}
    \label{fig:fusion_mlo_test}
\end{figure}

On retrouve, comme pour la plupart des architectures de ce mémoire, une
classe A totalement ignorée (0\,\% de rappel, seulement 10 échantillons de
test), et une confusion structurelle qui persiste entre les classes
adjacentes B/C et C/D.

\subsection{Fusion même-vue -- DenseNet121 + ResNet50 (vues CC)}
Même principe que la variante précédente, appliqué aux vues CC avec
DenseNet121 et ResNet50 comme paire de backbones.

\textbf{Terminé (10 époques).} Train Acc 90.01\,\%, Val Acc 75.69\,\%.
\textbf{Test Acc 76.35\,\%} (2000 images CC) : c'est le meilleur résultat
parmi les variantes de l'Approche 6 testées à ce jour, et il dépasse
légèrement la variante MLO ci-dessus.

\begin{figure}[H]
    \centering
    \includegraphics[width=0.8\textwidth]{image/courbe_accuracy_fusion_cc_densenet121_resnet50.png}
    \caption{DenseNet121+ResNet50 (fusion même-vue, CC) -- évolution de l'accuracy sur les 10 époques.}
    \label{fig:fusion_cc_courbe}
\end{figure}

Même constat que pour la variante MLO : la validation oscille, avec un
creux net à l'Époque 6 (60.68\,\%, en retrait par rapport aux 70.79\,\% de
l'époque précédente), avant de repartir à la hausse et de finir sur son
meilleur score à l'Époque 10. L'écart entre les deux courbes reste
cependant plus resserré que pour la variante MLO. La validation suit
d'assez près la tendance de l'entraînement à partir de l'Époque 7, ce qui
suggère une convergence un peu plus stable pour cette combinaison de
backbones sur les vues CC.

\begin{figure}[H]
    \centering
    \includegraphics[width=0.8\textwidth]{image/entrainement_fusion_cc_densenet121_resnet50.png}
    \caption{DenseNet121+ResNet50 (fusion même-vue, CC) -- résultat de l'entraînement (matrice de confusion sur validation, meilleure époque).}
    \label{fig:fusion_cc_entrainement}
\end{figure}

Ici, la classe A obtient sa meilleure précision de tout le mémoire sur un
jeu de validation (0.75, pour seulement 3 succès sur 8 échantillons). C'est
encourageant, mais il faut relativiser : avec un support aussi faible, un
seul échantillon de plus ou de moins suffirait à faire varier cette
précision de plusieurs dizaines de points. Le reste de la matrice suit le
même schéma que la variante MLO, avec l'essentiel des erreurs concentrées
entre classes voisines (B/C et C/D) plutôt que dispersées au hasard --
signe que le modèle capture bien une notion de densité progressive, même
s'il peine encore à tracer des frontières nettes entre catégories
adjacentes.

\begin{table}[h]
\centering
\caption{DenseNet121+ResNet50 (fusion même-vue, CC) -- Rapport de classification (test, 2000 images)}
\begin{tabular}{lrrrr}
\toprule
Classe & Précision & Rappel & F1-score & Support \\
\midrule
Density A & 0.50 & 0.10 & 0.17 & 10 \\
Density B & 0.55 & 0.86 & 0.67 & 190 \\
Density C & 0.95 & 0.74 & 0.83 & 1530 \\
Density D & 0.46 & 0.87 & 0.60 & 270 \\
\midrule
\textbf{Accuracy globale} & \multicolumn{3}{c}{} & \textbf{0.76} (2000) \\
Moyenne macro & 0.61 & 0.64 & 0.57 & \\
Moyenne pondérée & 0.84 & 0.76 & 0.78 & \\
\bottomrule
\end{tabular}
\end{table}

\begin{figure}[H]
    \centering
    \includegraphics[width=0.8\textwidth]{image/test._fusion_cc_densenet121_resnet50.png}
    \caption{DenseNet121+ResNet50 (fusion même-vue, CC) -- résultat du test final.}
    \label{fig:fusion_cc_test}
\end{figure}

Contrairement à la variante MLO, cette configuration parvient à reconnaître
correctement un échantillon de classe A (précision 0.50, rappel 0.10). Le
signal est faible mais non nul, à mettre en perspective avec l'échec
complet observé pour cette même classe sur toutes les autres architectures
de ce mémoire (voir notamment Section~\ref{sec:test-seuil-a}).

\subsection{Synthèse des architectures de l'Approche 6}
\begin{table}[h]
\centering
\caption{Comparaison des quatre variantes de l'Approche 6}
\begin{tabular}{lrrr}
\toprule
Variante & Train Acc & Val Acc & Test Acc \\
\midrule
Fusion même-vue MLO (RexNet150+ResNet50) & 88.74\,\% & 74.62\,\% & 74.85\,\% \\
Fusion même-vue CC (DenseNet121+ResNet50) & 90.01\,\% & 75.69\,\% & \textbf{76.35\,\%} \\
Architecture originale (RexNet150/CC+ResNet50/MLO) & 77.01\,\% & 85.80\,\% & 83.00\,\% \\
Variante siamoise (EfficientNet-B0, poids partagés) & 79.78\,\% & 84.38\,\% & \textbf{85.00\,\%} \\
\bottomrule
\end{tabular}
\end{table}

Sur les quatre variantes, la \textbf{variante siamoise obtient le meilleur
Test Acc du mémoire pour l'Approche 6} (85\,\%), suivie de près par
l'architecture originale (83\,\%). Les deux configurations à fusion même-vue
(MLO et CC), testées plus haut, restent en retrait (74.85\,\% et 76.35\,\%).
Le partage de poids entre vues (variante siamoise) semble donc apporter un
léger avantage par rapport aux branches hétérogènes indépendantes, malgré un
surapprentissage plus visible pendant l'entraînement.

% ============================================================================
\section{Approche 7 -- Classification hiérarchique en cascade (AB/CD $\to$ A/B \& C/D)}
\label{sec:approche7}

Cette approche reprend les deux mêmes backbones que l'Approche 2 (ViT et
ResNet50), mais remplace la classification directe à 4 sorties par une
décision en deux étapes, inspirée du raisonnement clinique par élimination.
Un premier classifieur binaire sépare les seins peu denses (A, B) des
seins denses (C, D). Un second classifieur, spécialisé, tranche ensuite à
l'intérieur du groupe retenu (A contre B, ou C contre D). Cette architecture
de classification hiérarchique en cascade est déjà présente dans le code du
projet (voir Section~\ref{sec:correspondance}).

\subsection{Résultats avec ViT}

\textbf{Terminé.} Test Acc \textbf{42.53\,\%} (4000 images), très en retrait
par rapport aux 70\,\% obtenus par l'Approche 2 avec le même backbone en
classification directe.

\begin{figure}[H]
    \centering
    \includegraphics[width=0.8\textwidth]{image/approche7_vit_test.png}
    \caption{Cascade hiérarchique (ViT) -- résultat du test final (4000 images), seuil de routage optimisé sur validation.}
    \label{fig:approche7_vit_test}
\end{figure}

La cause se trouve entièrement dans le routage de l'Étape 1. VinDr-Mammo
est dominé par les classes denses (C+D représentent 2240 des 4000 images de
test, contre 1760 pour A+B), et malgré une pondération de classe dans la
loss et un seuil optimisé sur une validation interne (0.5102, via critère de
Youden sur la courbe ROC), le classifieur binaire (A,B)/(C,D) reste biaisé
vers le groupe majoritaire.

\begin{figure}[H]
    \centering
    \includegraphics[width=0.55\textwidth]{image/approche7_vit_routage.png}
    \caption{Cascade hiérarchique (ViT) -- matrice de routage de l'Étape 1 (A,B vs C,D).}
    \label{fig:approche7_vit_routage}
\end{figure}

Les vrais cas C/D sont bien routés (rappel 86.96\,\%), mais plus de la
moitié des vrais cas A/B (926 sur 1760, soit 52.61\,\%) sont envoyés à tort
vers la branche C/D. Ce n'est pas un détail : une fois qu'un cas A ou B est
mal routé, il est perdu définitivement, puisque le Modèle 3 (C vs D) ne peut
plus jamais prédire A ou B. C'est exactement ce qu'on retrouve dans le
rapport de classification global -- le rappel de la classe B tombe à
9\,\%, et 350 des 800 vrais A (43.75\,\%) finissent classés C ou D.

Pour vérifier que ce n'est pas le Modèle 2 (A vs B) lui-même qui est en
cause, on peut l'évaluer isolément, indépendamment du routage de l'Étape 1
-- c'est-à-dire directement sur les 1760 images A+B du test, sans passer
par la décision préalable AB/CD.

\begin{figure}[H]
    \centering
    \includegraphics[width=0.55\textwidth]{image/approche7_vit_modele2.png}
    \caption{Cascade hiérarchique (ViT) -- Modèle 2 (A vs B) évalué isolément, hors effet du routage.}
    \label{fig:approche7_vit_modele2}
\end{figure}

Même isolé, le Modèle 2 ne dépasse que 52.90\,\% -- à peine mieux que le
hasard sur un problème à deux classes. Le vrai facteur limitant de cette
approche n'est donc pas seulement le routage : les trois modèles de la
cascade sont de petits MLP entraînés sur des features ViT déjà figées (sans
fine-tuning propre à chaque étage, contrairement à l'Approche 2 qui
fine-tune le backbone de bout en bout). La décomposition hiérarchique, ici,
amplifie l'effet du déséquilibre de classes au lieu de l'atténuer : une
erreur de routage précoce est irrécupérable, alors qu'un classifieur direct
à 4 sorties peut encore corriger le tir sur les étapes suivantes. Une
piste d'amélioration réaliste serait de fine-tuner le backbone séparément
pour chaque étage de la cascade plutôt que de réutiliser des features
gelées, mais cela représente un entraînement bien plus lourd que les MLP
utilisés ici.

\subsection{Résultats avec ResNet50}

\textbf{Terminé.} Test Acc \textbf{48.50\,\%} (4000 images) -- meilleur que
la version ViT (42.53\,\%), mais toujours très en retrait par rapport aux
74\,\% obtenus par l'Approche 2 avec le même backbone en classification
directe.

\begin{figure}[H]
    \centering
    \includegraphics[width=0.8\textwidth]{image/approche7_resnet50_test.png}
    \caption{Cascade hiérarchique (ResNet50) -- résultat du test final (4000 images), seuil de routage optimisé sur validation.}
    \label{fig:approche7_resnet50_test}
\end{figure}

Comme pour ViT, le routage de l'Étape 1 perd une part importante des cas :
sur les 4000 images de test, 974 (24.35\,\%) changent de groupe à tort
(890 vraies C/D routées vers A/B, 84 vraies A/B routées vers C/D), pour une
accuracy de routage globale de 75.65\,\%.

\begin{figure}[H]
    \centering
    \includegraphics[width=0.55\textwidth]{image/approche7_resnet50_routage.png}
    \caption{Cascade hiérarchique (ResNet50) -- matrice de routage de l'Étape 1 (A,B vs C,D).}
    \label{fig:approche7_resnet50_routage}
\end{figure}

Contrairement à la version ViT, où le maillon faible était le Modèle 2 (A
vs B, à peine 52.90\,\% isolé), c'est ici le \textbf{Modèle 3 (C vs D)} qui
pose problème. Le Modèle 2, évalué isolément sur les 400 vraies images A+B
sans passer par le routage, atteint \textbf{93.50\,\%} -- un score solide,
bien au-dessus du hasard.

\begin{figure}[H]
    \centering
    \includegraphics[width=0.55\textwidth]{image/approche7_resnet50_modele2.png}
    \caption{Cascade hiérarchique (ResNet50) -- Modèle 2 (A vs B) évalué isolément, hors effet du routage.}
    \label{fig:approche7_resnet50_modele2}
\end{figure}

Le problème se situe donc plus loin dans la cascade. Ce diagnostic ne
figure pas directement dans les logs, mais se déduit directement de la
matrice de confusion globale : une prédiction finale C ou D ne peut provenir
que d'un item correctement routé vers le groupe (C,D) à l'Étape 1 (le
groupe (A,B) ne produit jamais ce type de prédiction). En isolant les
colonnes C et D des lignes C et D de cette matrice, on obtient
$1394+928=2322$ vraies images C, et $132+256=388$ vraies images D,
correctement routées -- soit $2322+388=2710$ au total, ce qui recoupe
exactement la case correspondante de la matrice de routage de l'Étape 1
(Figure~\ref{fig:approche7_resnet50_routage}). Parmi ces 2710 images, le
Modèle 3 n'en classe correctement que $1394+256=1650$, soit
\textbf{60.9\,\%} -- à peine mieux que le hasard sur un problème binaire,
malgré une accuracy de 84.6\,\% atteinte pendant son propre entraînement
(30 époques, voir le détail des logs). Cet écart important entre l'accuracy
d'entraînement et la performance réelle en contexte de cascade est un signe
de surapprentissage du Modèle 3 sur ses propres données d'entraînement, qui
ne se traduit pas par une vraie capacité de généralisation sur la
distinction C/D. Combiné à la perte de routage, ce
double effet explique le score global modeste : sur les 3600 vraies images
C/D du test, seulement 1650 (45.83\,\%) ressortent correctement classées de
bout en bout de la cascade -- moins de la moitié.

Ce résultat, avec l'ordre différent des maillons faibles selon le backbone
(Modèle 2 pour ViT, Modèle 3 pour ResNet50), confirme le diagnostic déjà
posé pour ViT : la décomposition hiérarchique n'est pas fiable ici parce que
chaque étage repose sur des features gelées et un petit MLP, sans
fine-tuning propre à chaque étage. Le maillon le plus fragile varie d'un
backbone à l'autre, mais le résultat de bout en bout reste, dans les deux
cas, nettement inférieur à la classification directe à 4 sorties de
l'Approche 2 sur ce même backbone.

\section{Approche 8 -- Ensemble par fusion tardive (Approche 2 + Approche 3)}
\label{sec:approche8}

Les résultats des approches précédentes font apparaître un dilemme : l'Approche 3
(ResNet50 + Histogramme, vues MLO) atteint la meilleure accuracy globale de tout
le mémoire (85\,\%) mais ne reconnaît \emph{jamais} la classe A (rappel 0\,\%),
tandis que l'Approche 2 (ResNet50, fine-tuning direct) est nettement plus faible
en accuracy globale (74\,\%) mais reconnaît très bien cette même classe A
(rappel 90\,\%). Plutôt que de choisir arbitrairement entre les deux, on peut
tester s'il est possible de combiner leurs prédictions pour obtenir un
compromis meilleur que chacun des deux modèles pris isolément.

\subsection{Méthode}

La technique utilisée est un \textbf{ensemble par fusion tardive}, aussi appelé
\emph{vote souple} (\emph{soft voting}) ou \emph{late fusion} dans la littérature
sur l'apprentissage d'ensemble \cite{ref_dietterich}. Contrairement à une
architecture à double branche comme l'Approche 6 (Section~\ref{sec:approche6}),
où un unique modèle est entraîné de bout en bout sur les deux sources
d'information combinées (\emph{fusion précoce}), la fusion tardive combine les
\textbf{sorties} de deux modèles déjà entraînés \textbf{indépendamment}, sans
aucun réentraînement : chaque modèle garde son architecture et ses poids
d'origine, seule l'étape finale de décision est modifiée.

Pour une image $x$, soient $p^{(2)}(x) \in \mathbb{R}^4$ et $p^{(3)}(x) \in
\mathbb{R}^4$ les vecteurs de probabilités produits par la couche softmax de,
respectivement, l'Approche 2 et l'Approche 3. La probabilité d'ensemble est
obtenue par simple moyenne arithmétique :
\[
  p_{\text{ens}}(x) = \frac{1}{2}\Big(p^{(2)}(x) + p^{(3)}(x)\Big)
\]
et la classe prédite est celle qui maximise cette probabilité combinée :
\[
  \hat{y}_{\text{ens}}(x) = \operatorname*{arg\,max}_{c \,\in\, \{A,B,C,D\}} \; p_{\text{ens},c}(x)
\]
Il s'agit ici d'un cas particulier non pondéré ($w_1 = w_2 = 0.5$) de la
formulation générale de la fusion tardive à $K$ modèles,
$p_{\text{ens}}(x) = \sum_{k=1}^{K} w_k\, p^{(k)}(x)$ avec $\sum_k w_k = 1$ --
aucune pondération différenciée n'a été testée faute de temps, une piste
d'amélioration possible étant de calibrer $w_1, w_2$ sur un ensemble de
validation plutôt que de les fixer à $0.5$ chacun.

Comme l'Approche 3 ne fonctionne que sur les vues MLO, l'évaluation est faite
sur le sous-ensemble des 2000 images de test correspondant à cette vue (contre
4000 pour l'évaluation habituelle de l'Approche 2 seule, qui couvre aussi les
vues CC) -- les deux modèles sont donc comparés sur exactement le même
sous-ensemble d'images.

\subsection{Résultats}

\textbf{Terminé.} \textbf{Test Acc 80.55\,\%} (1611/2000 images MLO).

\begin{figure}[H]
    \centering
    \includegraphics[width=0.8\textwidth]{image/approche8_ensemble_test.png}
    \caption{Ensemble par fusion tardive (Approche 2 + Approche 3) -- résultat du test final (2000 images, vues MLO).}
    \label{fig:approche8_ensemble_test}
\end{figure}

\begin{table}[h]
\centering
\caption{Comparaison Approche 2 seule / Approche 3 seule / Ensemble (mêmes 2000 images de test, vues MLO)}
\begin{tabular}{lrrr}
\toprule
 & Approche 2 seule & Approche 3 seule & Ensemble \\
\midrule
Accuracy globale & 0.73 & 0.85 & 0.81 \\
Macro F1 & 0.60 & 0.54 & \textbf{0.62} \\
Rappel classe A & \textbf{0.90} (9/10) & 0.00 (0/10) & 0.30 (3/10) \\
\bottomrule
\end{tabular}
\end{table}

L'ensemble obtient le \textbf{meilleur macro F1 des trois} (0.62) : c'est la
seule des trois versions qui n'échoue catastrophiquement sur aucune classe.
Le rappel de la classe A passe de 0\,\% (Approche 3 seule) à 30\,\% -- un gain
réel, mais qui reste loin des 90\,\% obtenus par l'Approche 2 seule. La raison
est directe : la moyenne des probabilités est tirée vers le modèle le plus
confiant sur les gros volumes, à savoir l'Approche 3, qui a appris à ignorer
presque totalement la classe A. Sur les 10 vraies images A, l'ensemble en
reconnaît 3 et envoie les 7 autres vers la classe B (voisine en densité,
jamais vers C ou D) -- une erreur cohérente avec la proximité visuelle réelle
de A et B, pas un signe d'un modèle qui n'a rien appris.
L'accuracy globale (80.55\,\%) se situe, sans surprise, entre celle des deux
modèles pris isolément.

\section{Synthèse comparative : quel est le \og{}meilleur\fg{} modèle ?}
\label{sec:synthese-comparative}

Les résultats obtenus tout au long de ce chapitre montrent qu'il n'existe pas
un vainqueur unique parmi les approches testées, mais un arbitrage entre
trois critères qui ne pointent pas vers le même modèle. Le
Tableau~\ref{tab:synthese_comparative} résume cet arbitrage sur les
architectures les plus abouties du mémoire.

\begin{table}[h]
\centering
\caption{Synthèse comparative des approches selon trois critères}
\label{tab:synthese_comparative}
\begin{tabular}{lrrr}
\toprule
Approche & Test Acc & Rappel classe A & Macro F1 \\
\midrule
Approche 2 (ResNet50 seul) & 74{,}00\,\% & \textbf{0{,}95} & 0{,}62 \\
Approche 3 (ResNet50+Histogramme, MLO) & \textbf{85{,}00\,\%} & 0{,}00 & 0{,}54 \\
Approche 6 -- variante siamoise (EfficientNet-B0, CC+MLO) & \textbf{85{,}00\,\%} & -- & -- \\
Approche 8 -- Ensemble (Approche 2 + Approche 3) & 80{,}55\,\% & 0{,}30 & \textbf{0{,}62} \\
\bottomrule
\end{tabular}
\end{table}

\subsection{Critère 1 -- Accuracy globale}
Si l'on ne retient que l'exactitude globale, l'Approche 3 (ResNet50 +
Histogramme, vues MLO) et la variante siamoise de l'Approche 6 arrivent
en tête, à égalité, avec \textbf{85\,\%} de Test Acc -- les deux meilleurs
scores de tout le mémoire. C'est le critère le plus intuitif, mais aussi le
plus trompeur pris isolément : les deux modèles atteignent ce score en
sacrifiant presque entièrement la classe A (Section~\ref{sec:approche3}),
la classe minoritaire et cliniquement la plus délicate à repérer.

\subsection{Critère 2 -- Détection de la classe A}
Si l'on privilégie au contraire la détection de la classe A -- le tissu le
moins dense, dont la confusion avec les classes voisines a le plus
d'implications cliniques en dépistage -- l'Approche 2 seule (simple
fine-tuning direct de ResNet50, sans branche auxiliaire) domine largement
avec un rappel de \textbf{0{,}95} (19 images sur 20 correctement
identifiées), très loin devant toutes les architectures à double branche,
qui échouent toutes sur cette classe (rappel 0\,\%). C'est le résultat le
plus marquant du mémoire sur ce point précis (voir l'analyse détaillée,
Section~\ref{sec:approche3} ligne \og{}C'est le résultat le plus
marquant...\fg{}) : la simplicité de l'architecture semble ici aider le
réseau à apprendre un signal global (absence de texture dense sur
l'ensemble du cliché) que les architectures hybrides, focalisées sur la
fusion de caractéristiques locales, perdent en route.

\subsection{Critère 3 -- Équilibre entre les classes}
Si l'on cherche un compromis qui n'échoue catastrophiquement sur aucune des
quatre classes, l'Ensemble par fusion tardive (Approche 8) obtient le
meilleur macro F1 des trois versions comparées sur le même sous-ensemble
d'images (0{,}62, à égalité avec l'Approche 2 seule mais avec une accuracy
globale supérieure). Le rappel de la classe A y remonte à 0{,}30 (3/10),
un gain réel par rapport à l'Approche 3 seule (0{,}00) même s'il reste loin
des 0{,}90 de l'Approche 2 seule sur ce même sous-ensemble.

\subsection{Conclusion : pas de modèle \og{}final\fg{} unique}
Ce chapitre ne désigne donc pas un modèle \og{}gagnant\fg{} au sens
strict : le choix dépend de l'usage visé. Pour un système d'aide au
dépistage où manquer un cas de faible densité est plus coûteux qu'une
fausse alerte, l'Approche 2 seule reste la plus pertinente malgré son
accuracy globale plus modeste. Pour une estimation générale de la densité
sur l'ensemble de la population, l'Approche 3 ou la variante siamoise de
l'Approche 6 restent préférables. C'est cette absence de consensus qui
motive le choix du modèle utilisé comme point de départ pour
l'apprentissage fédéré au chapitre suivant : l'Approche 6 est
structurellement exclue (elle suppose une paire CC+MLO alignée côté client
H2, que le pipeline fédéré actuel ne construit pas, voir
Section~\ref{sec:federe}), tandis que les Approches 3, 4 et 5 seraient en
principe réutilisables (leur branche statistique se calcule sur une seule
image) mais demanderaient d'ajouter cette branche au modèle fédéré,
aujourd'hui un ResNet50 nu. Seule l'Approche 2, une architecture à une
seule branche déjà strictement identique au modèle fédéré existant, est
réutilisable sans aucune modification -- c'est elle qui est utilisée comme
point de départ pré-entraîné dans ce mémoire, faute de temps pour
implémenter les autres.

% ============================================================================
\chapter{Apprentissage fédéré (FedAvg)}
\label{sec:federe}

\section{Conception}
Conformément au sujet du mémoire, une simulation d'apprentissage fédéré a
été implémentée (\texttt{Codes/federated\_learning/}) suivant l'algorithme
FedAvg de McMahan et al.\ \cite{ref_fedavg}. L'idée de base est simple : au
lieu de rassembler toutes les images médicales à un seul endroit pour
entraîner un modèle, ce qui pose problème pour des données sensibles, on
envoie le \emph{modèle} chez chaque hôpital. Chaque hôpital l'entraîne chez
lui sur ses propres images, puis seuls les poids appris, jamais les images,
sont renvoyés à un serveur central qui les combine.

\begin{figure}[H]
\centering
\begin{tikzpicture}[
    scale=0.75, every node/.style={scale=0.75},
    node distance=1.3cm and 0.5cm,
    server/.style={rectangle, draw, fill=blue!10, text width=5.8cm, align=center, rounded corners, minimum height=1.1cm, font=\small\bfseries},
    client/.style={rectangle, draw, fill=green!10, text width=4.2cm, align=center, rounded corners, minimum height=1.1cm, font=\small},
    arrow/.style={-Stealth, thick, line width=1.2pt},
    labelstyle/.style={font=\scriptsize, color=black!80, align=center}
]
    \node (server) [server] {Serveur central \\ (simulation locale sur le cluster) \\ \textnormal{Agrégation : $W_{t+1} = \sum \frac{n_k}{n} W_t^k$}};
    \node (client1) [client, below left=of server] {Client H1 \\ VinDr-Mammo (public) \\ \textnormal{6801 images d'entraînement}};
    \node (client2) [client, below right=of server] {Client H2 \\ DDSM/CBIS-DDSM (public) \\ \textnormal{531 images d'entraînement}};
    \draw [arrow, color=blue!70!black] (server.west) -| node[labelstyle, left, pos=0.3] {1. Envoi du modèle \\ global ($W_t$)} (client1.north);
    \draw [arrow, color=blue!70!black] (server.east) -| node[labelstyle, right, pos=0.3] {1. Envoi du modèle \\ global ($W_t$)} (client2.north);
    \draw [arrow, color=orange!80!black] (client1.south) |- ++(0,-0.8) -| node[labelstyle, left, pos=0.7] {2. Renvoi des poids \\ entraînés ($W_t^1$)} (server.south);
    \draw [arrow, color=orange!80!black] (client2.south) |- ++(0,-0.8) -| node[labelstyle, right, pos=0.7] {2. Renvoi des poids \\ entraînés ($W_t^2$)} (server.south);
\end{tikzpicture}
\caption{Schéma du cycle d'apprentissage fédéré (FedAvg) tel qu'implémenté
dans ce projet, entre le serveur et les deux clients simulés H1 et H2.}
\label{fig:diag_apprentissage_federe}
\end{figure}

\textbf{Précision méthodologique importante.} Le terme \og{}simulation\fg{}
mérite d'être explicité : dans cette implémentation, l'ensemble du
protocole tourne sur une \textbf{seule machine physique} (un serveur GPU du
cluster de l'UMONS), pas sur deux serveurs distincts reliés par le réseau.
La \og{}fédération\fg{} entre H1 et H2 est obtenue par une simple boucle
qui traite d'abord les données de H1, puis celles de H2, l'une après
l'autre, dans le même processus Python -- ce même serveur joue donc
successivement les trois rôles du protocole (client H1, client H2, puis
serveur central qui agrège). Ce que le code appelle \og{}envoi\fg{} ou
\og{}renvoi\fg{} des poids correspond en réalité à des variables Python
passées en mémoire d'une fonction à l'autre, et non à une vraie
communication réseau entre machines (pas de \texttt{torch.distributed}, ni
de sockets, ni de connexion SSH entre serveurs). Le principe mathématique
de FedAvg (moyenne pondérée des poids de plusieurs clients, sans jamais
partager les images) est donc bien implémenté et respecté ; c'est
uniquement l'aspect \og{}deux hôpitaux physiquement séparés\fg{} qui est
simulé, faute de temps et d'accès à un second serveur dédié pour un
déploiement réellement distribué. Ce point a été discuté et validé comme
suffisant pour la démonstration attendue avec notre encadrant, M. Xavier
Lessage.

\textbf{Client H1} utilise VinDr-Mammo, le même dataset public que pour
toutes les autres approches du mémoire. \textbf{Client H2} devait à
l'origine représenter le réseau hospitalier privé HELORA, la cible réelle
du sujet de mémoire. Ces données n'étant pas accessibles pour cette
réplication, le dataset public CBIS-DDSM est utilisé à la place, comme
substitut technique. Ce choix n'est pas anodin : les deux datasets ont une
répartition de classes très différente (VinDr-Mammo est dominé par Density
C, alors que CBIS-DDSM est beaucoup plus équilibré, avec A=127, B=450,
C=465, D=329 sur 1371 images résolues). Cela recrée artificiellement une
situation proche de deux vrais hôpitaux avec des populations de patientes
différentes.

Le modèle utilisé est un simple ResNet50, une seule branche, sans
histogramme ni GLCM contrairement aux approches précédentes. Le protocole
FedAvg exige en effet que le modèle soit identique chez tous les clients ;
or le client H1 (\texttt{VinDrCCDataset}) est restreint aux vues CC par
construction, alors que le client H2 (\texttt{DDSMDataset}) ne filtre pas
par vue et mélange donc CC et MLO. Toutes les architectures explorées dans
ce mémoire ne sont pas concernées de la même façon par cette contrainte :
l'Approche 6 (et ses variantes) est structurellement exclue, puisqu'elle
suppose une paire CC+MLO alignée par sein, que le pipeline fédéré actuel
ne construit pas. Les Approches 3, 4 et 5, en revanche, n'ont pas besoin
d'une seconde vue -- leur branche histogramme/GLCM est calculée à partir de
la \emph{même} image que la branche CNN (voir
\texttt{HybridMammographyDataset.\_\_getitem\_\_}), pas d'une image
appariée. Rien n'empêche donc, en principe, de fédérer une architecture de
ce type, à condition (i) d'ajouter la branche statistique au modèle
fédéré, actuellement un ResNet50 nu, et (ii) de filtrer les deux clients
sur la même vue (par exemple MLO des deux côtés, ce que \texttt{DDSMDataset}
ne fait pas aujourd'hui). C'est une limite d'implémentation actuelle, pas
une impossibilité de principe -- une piste réaliste pour une itération
future de cette expérience plutôt que pour ce mémoire, faute de temps.

\section{Algorithme}
\begin{algorithm}[h]
\caption{FedAvg -- apprentissage fédéré entre les clients H1 et H2}
\begin{algorithmic}[1]
\Procedure{ServeurFedAvg}{}
  \State Initialiser les poids globaux $W_0$ (ResNet50, ImageNet)
  \For{chaque round $t = 1, \dots, T$}
    \For{chaque client $k \in \{\text{H1}, \text{H2}\}$}
      \State $W_t^k \gets$ \Call{ClientUpdate}{$k$, $W_{t-1}$}
    \EndFor
    \State $W_t \gets \sum_{k} \dfrac{n_k}{n}\, W_t^k$
    \Comment{agrégation pondérée par la taille locale $n_k$}
  \EndFor
  \State \Return $W_T$
\EndProcedure
\Statex
\Procedure{ClientUpdate}{$k$, $W$}
  \For{chaque époque locale $e = 1, \dots, E$}
    \For{chaque mini-lot $b$ du client $k$}
      \State $W \gets W - \eta \nabla_W \mathcal{L}(W; b)$
    \EndFor
  \EndFor
  \State \Return $W$ au serveur
\EndProcedure
\end{algorithmic}
\end{algorithm}

où $n_k$ est le nombre d'exemples d'entraînement du client $k$ et
$n = \sum_k n_k$. Dans notre implémentation, $E{=}1$ époque locale par round
est utilisée par défaut, ce qui rapproche le protocole d'une moyenne de
gradients quasi-synchrone plutôt que d'un fédéré à fort déséquilibre de
calcul local. Ce choix est motivé par la taille modeste des deux jeux de
données et le temps disponible.

\section{État d'avancement}
Le code est complet et a été testé localement : l'agrégation FedAvg a été
vérifiée mathématiquement (moyenne pondérée reproduite exactement sur un
cas de test synthétique), et le chargement des deux jeux de données a été
vérifié de bout en bout (lecture DICOM, prétraitement, batching).
L'entraînement a été lancé sur le cluster (GPU2 de deep2) le 4 août 2026, avec
$H1$ = 6801 images d'entraînement (VinDr-Mammo) et $H2$ = 531 images
(DDSM/CBIS-DDSM). Les 15 rounds prévus ont été menés à terme : chaque round
complet (2 \texttt{ClientUpdate} + agrégation + évaluation) a pris plusieurs
heures du fait de la lecture DICOM via NFS, voir contrainte de temps
discutée plus haut dans ce document.

\begin{table}[h]
\centering
\caption{Résultats FedAvg (E=1, 15 rounds) -- accuracy et loss de validation
par round}
\label{tab:federated_e3}
\begin{tabular}{lrrrrrr}
\toprule
Round & \multicolumn{3}{c}{Val Acc} & \multicolumn{3}{c}{Val Loss} \\
 & H1-VinDr & H2-DDSM & Global & H1-VinDr & H2-DDSM & Global \\
\midrule
1 & 65.33\% & 31.18\% & \textbf{62.86\%} & 0.7483 & 2.1213 & \textbf{0.8477} \\
2 & 60.33\% & 38.71\% & \textbf{58.77\%} & 0.8005 & 1.5975 & \textbf{0.8582} \\
3 & 70.08\% & 27.96\% & \textbf{67.03\%} & 0.6070 & 1.8421 & \textbf{0.6965} \\
4 & 75.25\% & 32.26\% & \textbf{72.14\%} & 0.5471 & 1.8859 & \textbf{0.6441} \\
5 & 59.75\% & 31.18\% & \textbf{57.68\%} & 0.9390 & 2.4917 & \textbf{1.0515} \\
6 & 74.33\% & 32.26\% & \textbf{71.29\%} & 0.5645 & 2.5631 & \textbf{0.7092} \\
7 & 61.33\% & 39.78\% & \textbf{59.77\%} & 0.8758 & 1.5269 & \textbf{0.9229} \\
8 & 62.50\% & 34.41\% & \textbf{60.47\%} & 0.8003 & 2.8950 & \textbf{0.9520} \\
9 & 64.42\% & 27.96\% & \textbf{61.78\%} & 0.7389 & 3.1998 & \textbf{0.9171} \\
10 & 71.00\% & 29.03\% & \textbf{67.96\%} & 0.5909 & 3.5098 & \textbf{0.8023} \\
11 & 70.25\% & 32.26\% & \textbf{67.50\%} & 0.6106 & 2.7237 & \textbf{0.7637} \\
12 & 69.92\% & 30.11\% & \textbf{67.03\%} & 0.5965 & 2.7623 & \textbf{0.7534} \\
13 & 69.92\% & 27.96\% & \textbf{66.88\%} & 0.6083 & 3.5827 & \textbf{0.8237} \\
14 & \textbf{78.67\%} & \textbf{41.94\%} & \textbf{76.01\%} & 0.4621 & 2.1410 & \textbf{0.5837} \\
15 & 70.83\% & 27.96\% & \textbf{67.73\%} & 0.6162 & 4.4000 & \textbf{0.8902} \\
\bottomrule
\end{tabular}
\end{table}

\begin{figure}[H]
    \centering
    \includegraphics[width=0.85\textwidth]{image/federated_rounds_accuracy.png}
    \caption{Évolution de l'accuracy de validation par round FedAvg (Global,
    H1-VinDr, H2-DDSM).}
    \label{fig:federated_rounds_accuracy}
\end{figure}

\begin{figure}[H]
    \centering
    \includegraphics[width=0.85\textwidth]{image/federated_rounds_loss.png}
    \caption{Évolution de la loss de validation par round FedAvg (Global,
    H1-VinDr, H2-DDSM).}
    \label{fig:federated_rounds_loss}
\end{figure}

L'écart entre les deux clients dès le premier round (65.33\,\% contre
31.18\,\%) est cohérent avec l'hétérogénéité non-IID volontairement
introduite entre les deux sources : $H2$ ne dispose que de 531 images
d'entraînement, contre 6801 pour $H1$, et présente une distribution de
densité différente, plus équilibrée que VinDr-Mammo (voir
Chapitre~\ref{sec:federe}). C'est précisément ce type d'écart que le
protocole FedAvg cherche à réduire au fil des rounds, en construisant un
modèle global qui profite de l'exposition aux deux distributions plutôt
qu'à une seule.

Le constat le plus net sur l'ensemble des 15 rounds n'est pas une
convergence progressive, mais un \textbf{verrouillage quasi total du score
global sur $H1$}. En comparant round par round la colonne \og{}Global\fg{}
à la colonne \og{}H1-VinDr\fg{} du Tableau~\ref{tab:federated_e3},
l'écart entre les deux reste remarquablement
constant, entre $-1.6$ et $-3.1$ points, à chacun des 15 rounds, quelle que
soit l'amplitude des variations par ailleurs (Round 1 : 65.33\,\% $\to$
62.86\,\%, écart $-2.47$ ; Round 4 : 75.25\,\% $\to$ 72.14\,\%, écart
$-3.11$ ; Round 14 : 78.67\,\% $\to$ 76.01\,\%, écart $-2.66$). Ce n'est pas
une coïncidence : la métrique \og{}Global\fg{} elle-même est une moyenne
pondérée $\text{acc}_{H1} \cdot \frac{n_{H1}}{n} + \text{acc}_{H2} \cdot
\frac{n_{H2}}{n}$, avec $n_{H1}/n \approx 93\,\%$. Le score global n'est
donc, par construction, presque rien d'autre qu'une version légèrement
abaissée du score $H1$ -- la contribution de $H2$ (7\,\% du poids) ne peut
mécaniquement pas déplacer le score global de plus de quelques points,
même quand $H2$ varie fortement (de 27.96\,\% à 41.94\,\% selon les
rounds).

Le même phénomène se retrouve dans l'agrégation des poids eux-mêmes ($W_{t+1}
= \sum_k \frac{n_k}{n} W_t^k$) : le modèle global ressemble
presque uniquement au modèle entraîné par $H1$ à chaque round, et n'intègre
que marginalement ce que $H2$ a appris localement. C'est la raison
structurelle pour laquelle $H2$ n'affiche, sur l'ensemble des 15 rounds,
aucune tendance claire à l'amélioration : son accuracy oscille sans
direction (31.18\,\%, 38.71\,\%, 27.96\,\%, 32.26\,\%, ..., jusqu'à un pic
isolé de 41.94\,\% au Round 14) plutôt que de converger, puisque le modèle
qu'il reçoit à chaque round est presque entièrement façonné par $H1$.

Le score $H1$ (et donc le score global) lui-même n'est pas stable : il
varie entre 59.75\,\% (Round 5) et 78.67\,\% (Round 14), sans palier net.
Avec $E{=}1$ seule époque locale par round, le modèle repart d'une remise à
zéro locale à chaque round puis n'a qu'une seule passe sur les données
avant d'être renvoyé et remoyenné -- ce qui laisse peu de temps pour
converger proprement d'un round à l'autre et explique cette instabilité
(par exemple Round 4 à 75.25\,\% suivi d'une chute au Round 5 à 59.75\,\%).
Le \textbf{Round 14 est le meilleur de toute la simulation} (H1 78.67\,\%,
H2 41.94\,\%, Global 76.01\,\%) et c'est donc son modèle
(\texttt{federated\_global\_best.pth}) qui est retenu pour l'évaluation
finale sur le jeu de test (Section~\ref{sec:federe-baseline} ci-dessous).

Ce résultat pointe vers deux limites structurelles de cette première
configuration (E=1, pas de plancher de poids pour $H2$) plutôt que vers un
simple manque de rounds : une seule époque locale ne laisse pas au modèle
le temps de converger avant chaque agrégation, et le poids proportionnel
au nombre d'images ($n_k/n$) rend le score global -- et le modèle lui-même
-- presque insensible à ce que $H2$ apprend. Ce sont exactement les deux
pistes testées dans le run suivant (E=6 époques locales, et un poids
plancher de 25\,\% garanti à $H2$ dans l'agrégation, voir
Section~\ref{sec:federe-e6} ci-dessous).

\section{Comparaison avec une baseline centralisée}
\label{sec:federe-baseline}

Pour juger si la contrainte de fédération coûte réellement de la
performance, une baseline centralisée a été entraînée : les deux jeux de
données ($H1$ = VinDr-Mammo, $H2$ = DDSM/CBIS-DDSM) sont regroupés dans un
seul entraînement classique (5 époques), sans aucune séparation entre
clients ni agrégation. Le modèle FedAvg final (meilleur round, le Round 14
du Tableau~\ref{tab:federated_e3}) a lui aussi été évalué sur les mêmes
jeux de \emph{test} réservés de chaque source, ce qui permet une
comparaison directe, homogène, plutôt qu'un rapprochement approximatif
entre une accuracy de test et une accuracy de validation.

\begin{table}[h]
\centering
\caption{Baseline centralisée vs. FedAvg (E=1, 15 rounds) -- résultats sur les jeux de test}
\label{tab:federated_vs_centralized}
\begin{tabular}{lrrr}
\toprule
Source & Baseline centralisée & FedAvg (E=1, Round 14) & Support \\
\midrule
VinDr-Mammo & \textbf{82.30\,\%} & 81.60\,\% & 2000 images \\
DDSM/CBIS-DDSM & 31.17\,\% & \textbf{38.31\,\%} & 154 images \\
\bottomrule
\end{tabular}
\end{table}

Le résultat est plus favorable à la fédération que ce que suggérait
l'ancienne comparaison (validation vs. test) : sur VinDr, l'écart entre
FedAvg et la baseline centralisée n'est que de 0.70 point (82.30\,\% contre
81.60\,\%) -- un coût de fédération quasi négligeable dans cette
configuration, bien loin de ce qu'on pourrait craindre a priori. Sur DDSM,
FedAvg fait même \textbf{mieux} que le centralisé (38.31\,\% contre
31.17\,\%, soit +7.14 points). Ce dernier point mérite d'être nuancé : le
jeu de test DDSM ne compte que 154 images, donc cet écart doit être lu avec
prudence plutôt que comme une preuve définitive que la fédération
\og{}aide\fg{} structurellement DDSM ; il montre au minimum que la
fédération n'impose pas ici de coût de performance sur ce client, contrairement à
l'hypothèse initiale.

Cette comparaison répond directement à la question de notre encadrant
(\og{}est-ce que ça a permis d'avoir de meilleurs résultats ?\fg{}) : dans
cette première configuration (E=1, sans poids plancher), FedAvg n'apporte
pas de gain net mesurable sur VinDr par rapport au centralisé, mais ne
coûte pas non plus grand-chose -- et fait mieux sur DDSM. La vraie valeur
de l'apprentissage fédéré ici n'est donc pas un gain de performance brute,
mais la capacité à entraîner un modèle exploitant les deux sources
\emph{sans jamais centraliser les images}, ce qui reste le problème que
FedAvg est censé résoudre en premier lieu (voir la motivation de
confidentialité, Section~\ref{sec:etat-art-federe}).

\section{Point de départ pré-entraîné (Approche 2) plutôt qu'ImageNet}
\label{sec:federe-pretrained}

Notre encadrant a insisté sur un point méthodologique important : avant de
faire de l'apprentissage fédéré, il faut déjà disposer d'un modèle qui
fonctionne bien, et tester si possible FedAvg à partir de ce meilleur
modèle plutôt que de partir de zéro. Comme discuté en
Section~\ref{sec:synthese-comparative}, aucune de nos meilleures approches
(Approche 3, Approche 6) n'est directement réutilisable comme socle fédéré
-- seule l'Approche 2 (ResNet50, une seule branche) partage l'architecture
exigée par le protocole fédéré actuel. Un second run a donc été lancé,
strictement identique au premier (E=1, 15 rounds, pas de poids plancher),
mais en initialisant $W_0$ avec les poids déjà fine-tunés de l'Approche 2
sur VinDr-Mammo (74\,\% Test Acc en seul, Section~\ref{sec:approche2})
plutôt qu'avec des poids ImageNet bruts.

\begin{table}[h]
\centering
\caption{FedAvg (E=1, 15 rounds) -- ImageNet vs. poids pré-entraînés
(Approche 2), résultats sur le jeu de test}
\label{tab:federated_pretrained}
\begin{tabular}{lrrrr}
\toprule
 & VinDr Test Acc & VinDr Macro F1 & VinDr Rappel A & DDSM Test Acc \\
\midrule
$W_0$ = ImageNet & 81.60\,\% & 0.57 & 10\,\% (1/10) & \textbf{38.31\,\%} \\
$W_0$ = Approche 2 & \textbf{81.80\,\%} & \textbf{0.66} & \textbf{60\,\% (6/10)} & 32.47\,\% \\
\bottomrule
\end{tabular}
\end{table}

Le résultat est un vrai compromis, pas un gain net dans un sens ou dans
l'autre. Sur VinDr, l'accuracy globale ne bouge presque pas (81.60\,\%
$\to$ 81.80\,\%), mais l'\textbf{équilibre entre classes s'améliore
nettement} : le macro F1 passe de 0.57 à 0.66, porté surtout par un rappel
de la classe A qui bondit de 10\,\% (1/10) à 60\,\% (6/10) -- cohérent avec
le fait que l'Approche 2 est justement, dans tout ce mémoire, le modèle le
plus fort sur cette classe précise (rappel 0.95 en seul,
Section~\ref{sec:approche2}) ; une partie de cette force semble transférée
au modèle fédéré malgré les 15 rounds d'agrégation qui suivent. En
revanche, sur DDSM, c'est une régression : 38.31\,\% $\to$ 32.47\,\%. Partir
de poids entraînés uniquement sur VinDr n'aide pas -- et gêne même
légèrement -- la généralisation vers DDSM, dont la distribution d'image
est différente (mammographie numérique native contre films argentiques
numérisés, voir Chapitre~\ref{sec:federe}).

Ce résultat illustre bien pourquoi il ne faut pas s'attendre à un gain
uniforme : partir du \og{}meilleur\fg{} modèle compatible améliore la
robustesse du modèle fédéré sur le client dont il est issu (VinDr, et
spécifiquement sa classe la plus difficile), sans pour autant garantir de
meilleurs résultats sur l'autre client. Avec 154 images de test côté DDSM
(dont seulement 15 en classe A côté VinDr), ces écarts doivent être lus
avec la prudence habituelle sur de petits effectifs, mais la direction du
compromis (VinDr mieux équilibré, DDSM moins bon) est cohérente entre
validation et test et ne semble donc pas être un simple artefact de bruit.

\subsection{Comparaison directe avec l'Approche 2 seule, à jeu de test identique}

L'Approche 2 seule affiche 74\,\% de Test Acc (Section~\ref{sec:approche2}).
Ce chiffre porte sur 4000 images, CC et MLO mélangées. Le modèle fédéré, lui,
n'entraîne et n'évalue que sur des vues CC -- 2000 images. Comparer 74\,\%
à 81.60\,\% ou 81.80\,\% reviendrait donc à comparer deux épreuves
différentes.

Le checkpoint de l'Approche 2 a été réévalué sur ces mêmes 2000 images CC,
rien de plus qu'une passe d'inférence sur le modèle déjà entraîné. Le score
tombe à \textbf{75.40\,\%}. C'est ce chiffre qu'il faut opposer aux résultats
fédérés, pas le 74\,\% d'origine.

\begin{table}[h]
\centering
\caption{Test Acc sur le même jeu de 2000 images CC -- Approche 2 seule vs. FedAvg}
\label{tab:federated_vs_approche2_cc}
\begin{tabular}{lr}
\toprule
Modèle & Test Acc (2000 img CC) \\
\midrule
Approche 2 seule (sans fédération) & 75.40\,\% \\
FedAvg E=3, $W_0$ = ImageNet & 81.60\,\% \\
FedAvg E=3, $W_0$ = Approche 2 & \textbf{81.80\,\%} \\
\bottomrule
\end{tabular}
\end{table}

Sur ce jeu de test identique, FedAvg dépasse l'Approche 2 seule de 6.20 à
6.40 points. C'est la réponse à la question posée par notre encadrant :
oui, la fédération fait mieux que le meilleur modèle compatible entraîné
seul, sur cette même donnée de test.

Le rappel de la classe A grimpe à 100\,\% (10/10) pour l'Approche 2 seule
sur ce sous-ensemble CC, contre 95\,\% sur les 4000 images complètes. Dix
images, c'est peu. Un score parfait sur un effectif aussi petit se lit avec
prudence, mais confirme la même tendance observée partout ailleurs dans ce
mémoire : cette architecture reconnaît bien la classe A.

\section{Davantage d'époques locales et poids plancher pour H2 (E=6, 17 rounds)}
\label{sec:federe-e6}

Un troisième run, combinant les deux autres pistes identifiées plus haut
(Section~\ref{sec:federe-baseline}) -- davantage d'époques locales par
round ($E{=}6$ au lieu de $E{=}1$) et un poids plancher de 25\,\% garanti à
$H2$ dans l'agrégation (\texttt{min\_weight=0.25}), indépendamment de son
nombre réel d'images -- a été lancé sur 17 rounds. \todo{Run en cours au
moment de la rédaction (Round 9/17 atteint) -- compléter cette section avec
le tableau round par round et l'évaluation finale sur le jeu de test une
fois les 17 rounds terminés.}

% ============================================================================
\chapter{Interprétabilité et visualisation du modèle}
\label{sec:interpretabilite}

Lors du point d'avancement avec notre encadrant, celui-ci a insisté sur un
point : un chiffre d'accuracy seul ne suffit pas à convaincre un jury. Il
faut pouvoir montrer que l'espace de représentation appris par le réseau
organise réellement les quatre classes de densité de façon cohérente, et
que les prédictions s'appuient sur des zones anatomiquement pertinentes du
cliché plutôt que sur des artefacts -- fond noir, pli de peau, ou pire, la
présence d'une anomalie. Deux techniques complémentaires répondent à cette
double exigence : une visualisation t-SNE de l'espace latent, pour la
structure globale sur l'ensemble du jeu de test, et une analyse Grad-CAM au
niveau pixel, pour une explication locale image par image.

Le modèle retenu pour cette analyse est l'Approche 3 (ResNet50 +
Histogramme, vues MLO) : c'est le meilleur résultat obtenu sur l'ensemble
du mémoire (Test Acc 85\,\%, Section~\ref{sec:approche3}), et son backbone
CNN se prête naturellement à Grad-CAM, qui a besoin d'une carte
d'activation spatiale issue d'une couche convolutionnelle -- contrairement
à un ViT, où la même analyse demanderait une méthode adaptée à
l'attention plutôt qu'à la convolution.

\section{Visualisation de l'espace latent par réduction de dimension (t-SNE)}
\label{sec:tsne}

Le vecteur de caractéristiques produit par la branche image, juste avant la
couche de classification, vit dans un espace de plusieurs centaines de
dimensions -- impossible à représenter tel quel. La technique t-SNE
(\emph{t-distributed Stochastic Neighbor Embedding}) \cite{ref_tsne}
projette ce nuage de points de haute dimension dans un plan 2D, en cherchant
à préserver au mieux les relations de voisinage entre points proches dans
l'espace d'origine. Concrètement, pour chaque image du jeu de test, on
extrait le vecteur produit par l'avant-dernière couche du réseau via un
\emph{hook} PyTorch placé sur la couche cible (voir
\texttt{Codes/interpretability\_analysis.py}). L'ensemble de ces vecteurs
est ensuite projeté en 2D avec la bibliothèque scikit-learn.

\begin{figure}[H]
    \centering
    \includegraphics[width=0.85\textwidth]{image/tsne_approche3_resnet50_mlo.png}
    \caption{Visualisation t-SNE des caractéristiques apprises par le modèle
    retenu (Approche 3, ResNet50 + Histogramme, MLO). Chaque point représente
    une mammographie du jeu de test projetée dans l'espace latent 2D,
    colorée selon sa classe de densité réelle (300 images échantillonnées
    aléatoirement, seed=42).}
    \label{fig:tsne}
\end{figure}

Si le modèle a effectivement appris une représentation utile de la densité,
on s'attend à une organisation continue de l'espace projeté : les
échantillons de la classe A (tissu adipeux) devraient occuper une zone
plutôt séparée des échantillons D (tissu extrêmement dense), avec les
classes B et C dans une position intermédiaire. Il n'y a pas forcément de
frontière nette entre elles, puisque la densité mammaire est par nature un
continuum plutôt que quatre catégories parfaitement disjointes. À
l'inverse, un nuage de points totalement mélangé, sans structure
identifiable, indiquerait que l'espace de représentation appris par le
réseau ne capture pas l'information de densité.

Dans les faits, la projection montre une organisation partielle plutôt
qu'un mélange complet. La classe B forme un groupe assez net et bien
séparé du reste, en haut de la projection. La classe D occupe
majoritairement une zone distincte en bas à gauche, mais plusieurs de ses
points se retrouvent isolés loin à l'intérieur de la zone dominée par C --
cohérent avec la confusion C/D déjà observée dans les matrices de
confusion de ce mémoire. La classe C, très majoritaire dans le jeu de
données, occupe naturellement une large portion de l'espace.

La classe A, elle, n'apparaît qu'à travers un unique point sur toute la
projection. Ce n'est pas un défaut du modèle mais un effet direct de sa
rareté : avec seulement une dizaine d'images de cette classe dans les 2000
images de test, un échantillon aléatoire de 300 images n'en retient
presque jamais plus d'une ou deux. Il est donc impossible de tirer une
conclusion sur la façon dont le modèle organise cette classe à partir de
cette seule figure -- c'est une limite liée à la disponibilité des
données, pas une observation à interpréter.

\subsection{Comparaison complémentaire : train vs test}

En complément de la figure précédente, le même t-SNE a été ajusté
conjointement sur un échantillon du \textbf{train} et un échantillon du
\textbf{test} (300 images de chaque, projetées dans le même espace 2D),
pour vérifier si le modèle organise l'espace latent de la même façon sur
les deux, ou si le test forme des zones à part -- signe d'un décalage de
distribution ou d'un surapprentissage sur le train.

\begin{figure}[H]
    \centering
    \includegraphics[width=0.85\textwidth]{image/tsne_approche3_resnet50_mlo_traintest.png}
    \caption{Même visualisation t-SNE que la Figure~\ref{fig:tsne}, avec en
    plus un échantillon du train projeté dans le même espace 2D. La couleur
    encode toujours la classe de densité réelle ; la transparence encode le
    split (train = clair, test = foncé) -- 300 images de chaque split,
    seed=42.}
    \label{fig:tsne_traintest}
\end{figure}

Le point le plus rassurant de cette figure est que, à l'intérieur d'une
même classe, les points clairs (train) et foncés (test) occupent
globalement les \textbf{mêmes régions} de l'espace projeté, sans zone de
test isolée à l'écart du nuage d'entraînement. C'est le signe attendu d'un
modèle qui généralise correctement -- s'il avait mémorisé le train plutôt
qu'appris une notion générale de densité, on s'attendrait à voir les deux
splits former des amas séparés, même à l'intérieur d'une classe donnée.

\section{Analyse de l'interprétabilité clinique (Grad-CAM)}
\label{sec:gradcam}

\subsection{Motivation}

Un réseau de neurones profond reste, par construction, une boîte noire : on
observe la prédiction en sortie sans savoir sur quelles zones de l'image
elle s'appuie réellement. Pour un usage clinique, cette opacité pose un
vrai problème de confiance -- un radiologue ne peut pas se fier à un outil
dont il ignore s'il se base sur le tissu mammaire ou sur un artefact de
numérisation. C'est exactement le risque soulevé par notre encadrant : si
le modèle apprenait à estimer la densité en s'appuyant en partie sur la
présence d'une anomalie, bénigne ou maligne, plutôt que sur la texture du
parenchyme sain, la mesure de densité perdrait tout son sens clinique,
puisque densité et détection d'anomalie sont deux tâches bien distinctes.

L'algorithme Grad-CAM (\emph{Gradient-weighted Class Activation Mapping}),
proposé par Selvaraju \textit{et al.} \cite{ref_gradcam}, répond à cette
question en générant une carte de chaleur superposée à l'image d'origine,
qui indique les régions ayant le plus contribué à la décision du réseau
pour une classe donnée. Cette approche a déjà été appliquée avec succès à
l'analyse de mammographies dans plusieurs travaux récents
\cite{ref_gradcam_mammo1, ref_gradcam_mammo2}, ce qui conforte le choix de
cette méthode plutôt que, par exemple, une approche par occlusion,
beaucoup plus coûteuse en temps de calcul.

\subsection{Formulation mathématique}

Pour une classe cible $c$, soit $A^k$ la carte d'activation du $k$-ième
filtre de la dernière couche convolutive et $y^c$ le score (le logit, avant
la fonction Softmax) attribué par le réseau à la classe $c$. L'importance de
chaque filtre pour cette classe s'obtient en moyennant spatialement le
gradient du score par rapport à la carte d'activation correspondante :
\begin{equation}
    \alpha_k^c = \frac{1}{Z} \sum_{i=1}^{H} \sum_{j=1}^{W} \frac{\partial y^c}{\partial A_{i,j}^k}
\end{equation}
où $Z = H \times W$ est la dimension spatiale de la carte d'activation. La
carte de chaleur finale est la combinaison linéaire des cartes d'activation
pondérées par ces coefficients, suivie d'une ReLU pour ne conserver que les
contributions qui renforcent la classe cible :
\begin{equation}
    L_{\text{Grad-CAM}}^c = \text{ReLU}\left( \sum_k \alpha_k^c A^k \right)
\end{equation}
Cette carte a la taille de la dernière couche convolutive, typiquement
$7 \times 7$ pour les backbones utilisés dans ce projet. Elle est ensuite
ré-échantillonnée par interpolation bilinéaire à la taille de l'image
d'entrée, pour être superposée au cliché original. L'implémentation
utilisée ici (\texttt{Codes/interpretability\_analysis.py}, classe
\texttt{GradCAM}) repose sur les \emph{hooks} PyTorch, un mécanisme qui
permet d'intercepter les activations et les gradients d'une couche donnée
pendant les passes avant et arrière du réseau, sans modifier l'architecture
du modèle.

\begin{figure}[H]
    \centering
    \includegraphics[width=0.98\textwidth]{image/gradcam_approche3_resnet50_mlo.png}
    \caption{Grad-CAM sur le modèle retenu (Approche 3, ResNet50 +
    Histogramme, MLO) : trois exemples par classe de densité (au lieu d'un
    seul), pour vérifier la stabilité du comportement du modèle plutôt que
    de se fier à un seul cas potentiellement non représentatif. Pour
    chaque exemple, l'image originale est suivie de sa carte de chaleur
    Grad-CAM correspondante. Colorbar verticale commune à toute la figure.}
    \label{fig:gradcam}
\end{figure}

Précision méthodologique importante pour la lecture de cette figure : pour
chaque image, le calcul du gradient est forcé sur la \emph{vraie} classe de
l'image (et non sur la classe que le modèle aurait spontanément prédite).
Le pourcentage affiché dans chaque sous-titre est donc la probabilité que
le modèle attribue à la bonne réponse, pas nécessairement sa prédiction la
plus probable -- une nuance essentielle pour les images de la classe A
ci-dessous.

\subsection{Points de vigilance à vérifier sur le modèle final}

Notre encadrant avait relevé deux éléments lors d'une analyse Grad-CAM
préliminaire menée sur une architecture antérieure (RexNet-150 sur vues
MLO, en dehors du présent projet). Il faudra les vérifier à nouveau une
fois la même analyse reproduite sur le modèle actuel :
\begin{itemize}
    \item pour la classe A, la carte de chaleur tendait par moments à
    réagir sur le fond noir de l'image ou sur un pli de peau plutôt que sur
    le tissu mammaire lui-même -- un recadrage plus strict avant le passage
    dans le réseau pourrait limiter ce biais ;
    \item il faut s'assurer que le modèle ne s'appuie pas sur la présence
    d'une anomalie (bénigne ou maligne) pour estimer la densité. Une
    vérification possible, si le temps le permet, consisterait à évaluer le
    modèle spécifiquement sur un sous-ensemble d'images classées BI-RADS 1-2
    (négatives) et à comparer les cartes obtenues à celles d'images
    contenant une anomalie confirmée.
\end{itemize}

Sur le premier point, la figure obtenue est plutôt rassurante : pour les
quatre classes, les zones d'activation (rouge/jaune) restent concentrées
sur le tissu mammaire lui-même, en particulier autour du mamelon et des
régions les plus claires du cliché -- aucune des quatre cartes ne réagit
de façon marquée sur le fond noir. Le second point, en revanche, demanderait
la comparaison systématique proposée ci-dessus (BI-RADS négatif vs positif)
pour être tranché avec certitude ; faute de temps, il reste une piste à
explorer plutôt qu'une conclusion établie ici.

Classe par classe, sur les trois exemples testés : la classe \textbf{A}
présente systématiquement une activation étalée sur une large portion du
sein plutôt qu'un foyer localisé, et surtout une probabilité assignée à la
bonne classe très basse sur les trois cas (4.8\,\%, 6.1\,\%, 2.7\,\%) -- le
modèle est ici franchement et \emph{constamment} dans l'erreur, cohérent
avec le rappel quasi nul déjà observé sur cette classe tout au long du
mémoire (Section~\ref{sec:test-seuil-a}). La classe \textbf{C} montre un
foyer d'activation localisé et très stable, centré sur la masse de tissu
dense près du mamelon, avec une probabilité correcte quasiment identique
sur les trois exemples (48.7\,\%, 49.3\,\%, 47.8\,\%). La classe \textbf{B}
présente elle aussi un foyer localisé, avec une probabilité correcte plus
variable (48.3\,\%, 43.8\,\%, 78.8\,\%), mais toujours supérieure à 40\,\%.

La classe \textbf{D} est la plus instructive des quatre une fois plusieurs
exemples comparés : les deux premiers cas donnent un foyer concentré sur la
zone dense avec une probabilité correcte raisonnable (39.4\,\% et
42.9\,\%), mais le troisième s'effondre à seulement 2.3\,\% -- une
instabilité totalement invisible sur la version à un seul exemple par
classe utilisée précédemment. C'est exactement l'intérêt d'étendre cette
analyse à plusieurs exemples : un seul cas, choisi par hasard, aurait pu
donner une impression de fiabilité pour la classe D, alors qu'un exemple
sur trois révèle une erreur franche. Dans l'ensemble, le modèle s'appuie
sur les bonnes structures anatomiques (le tissu fibro-glandulaire, plus
clair sur le cliché) pour les classes B, C et D, avec une localisation de
l'activation cohérente d'un exemple à l'autre -- mais son niveau de
confiance sur D n'est pas garanti d'une image à l'autre, contrairement à C.
La classe A, elle, reste peu fiable de façon constante -- ce qui confirme,
au niveau visuel cette fois, un problème déjà identifié statistiquement
partout ailleurs dans ce mémoire : le manque de données pour cette classe
limite sa capacité d'apprentissage, indépendamment de l'architecture
utilisée.

\vspace{0.3cm}
\noindent\textit{Note pour la bibliographie (\texttt{.bib}) :} les clés
utilisées ci-dessus correspondent aux références suivantes, à ajouter si
elles n'existent pas déjà dans le fichier de bibliographie du mémoire :
\texttt{ref\_gradcam} = Selvaraju et al., \emph{Grad-CAM: Visual Explanations
from Deep Networks via Gradient-based Localization}, ICCV 2017
(arXiv:1610.02391) ; \texttt{ref\_tsne} = van der Maaten \& Hinton,
\emph{Visualizing Data using t-SNE}, JMLR 2008 ;
\texttt{ref\_gradcam\_mammo1} = Al Mansour et al. (ou titre exact de l'article
EfficientNetV2 + Grad-CAM++), Diagnostics 2026, \url{mdpi.com/2075-4418/16/1/105} ;
\texttt{ref\_gradcam\_mammo2} = Talaat et al., \emph{Grad-CAM Enabled Breast
Cancer Classification with a 3D Inception-ResNet V2}, Cancers 2024,
\url{mdpi.com/2072-6694/16/21/3668}.

\section{Précision et rappel de la classe A en fonction du seuil de décision}
\label{sec:seuil-classe-a}

Les métriques présentées jusqu'ici pour la classe A (précision, rappel)
sont calculées avec la règle de décision par défaut : la classe prédite est
celle de plus haute probabilité (seuil implicite de 0.5 dans un contexte
binaire, ou simplement l'argmax en multi-classe). Cette règle n'est pas
forcément optimale pour une classe aussi minoritaire -- il est instructif
de regarder comment précision et rappel évoluent si l'on fait varier
explicitement le seuil de décision sur $P(\text{classe}=A)$, plutôt que de
se limiter à ce seul point de fonctionnement. Cette analyse est faite à la
fois sur l'Approche 2 (ResNet50, fine-tuning direct) et l'Approche 3
(ResNet50 + Histogramme, MLO), en réutilisant les probabilités déjà
calculées pour la fusion tardive de l'Approche 8
(Section~\ref{sec:approche8}), sur le même sous-ensemble de 2000 images de
test (vues MLO).

\begin{figure}[H]
    \centering
    \includegraphics[width=0.95\textwidth]{image/courbe_seuil_classeA.png}
    \caption{Précision et rappel de la classe A en fonction du seuil de
    décision sur $P(\text{classe}=A)$, pour l'Approche 2 (gauche) et
    l'Approche 3 (droite), sur les mêmes 2000 images de test (vues MLO).}
    \label{fig:seuil_classeA}
\end{figure}

Le contraste entre les deux courbes est net. Pour l'\textbf{Approche 3}
(droite), le rappel s'effondre à 0 dès un seuil d'environ 0.10 : le modèle
n'assigne pratiquement jamais plus de 10\,\% de probabilité à la classe A,
même sur les vraies images de cette classe -- ce n'est donc pas un problème
de seuillage mal choisi, mais un problème d'apprentissage : le signal
distinctif de la classe A n'a simplement pas été capturé par ce modèle,
quel que soit le seuil retenu. Pour l'\textbf{Approche 2} (gauche), le
rappel reste élevé (autour de 90\,\%) jusqu'à un seuil proche de 0.75-0.80,
avec une précision qui progresse régulièrement de pair -- une vraie zone de
compromis exploitable existe ici, contrairement à l'Approche 3. Ce résultat
confirme visuellement, sur l'ensemble du spectre des seuils possibles et
pas seulement au point de fonctionnement par défaut, que la classe A est
un signal que l'Approche 2 a réellement appris à reconnaître, là où
l'Approche 3 l'a essentiellement ignorée pendant l'entraînement.

\section{Analyse comparative -- Approche 2 (préparée en parallèle)}
\label{sec:interpretabilite-approche2}

Le choix du modèle final entre l'Approche 2 et l'Approche 3
(Section~\ref{sec:approche8}) est en attente de l'avis de notre encadrant,
compte tenu du compromis accuracy globale / rappel de la classe A détaillé
plus haut. Par précaution, la même analyse t-SNE + Grad-CAM que celle menée
sur l'Approche 3 a donc été reproduite à l'identique sur l'Approche 2
(ResNet50, fine-tuning direct), afin d'avoir ce chapitre prêt quelle que
soit la réponse -- l'une des deux versions sera retirée du mémoire final
une fois la décision connue.

\begin{figure}[H]
    \centering
    \includegraphics[width=0.85\textwidth]{image/tsne_approche2_resnet50.png}
    \caption{Visualisation t-SNE des caractéristiques apprises par
    l'Approche 2 (ResNet50, fine-tuning direct), jeu de test seul, même
    protocole que la Figure~\ref{fig:tsne}.}
    \label{fig:tsne_approche2}
\end{figure}

\begin{figure}[H]
    \centering
    \includegraphics[width=0.85\textwidth]{image/tsne_approche2_resnet50_traintest.png}
    \caption{Même visualisation t-SNE que la Figure~\ref{fig:tsne_approche2},
    avec en plus un échantillon du train projeté dans le même espace 2D
    (transparence = split, train clair / test foncé), même protocole que la
    Figure~\ref{fig:tsne_traintest}.}
    \label{fig:tsne_approche2_traintest}
\end{figure}

La structure obtenue est comparable à celle de l'Approche 3 : les classes
B et D forment chacune une zone assez distincte, la classe C domine
largement l'espace, et train et test se superposent globalement bien à
l'intérieur de chaque classe -- même absence de séparation train/test que
pour l'Approche 3, donc pas de signe de surapprentissage plus marqué de ce
côté. Un point isolé de la classe A est visible ici (contre un seul
également pour l'Approche 3), toujours trop peu nombreux pour en tirer une
conclusion sur son organisation dans l'espace latent. On note aussi un
petit amas isolé, loin du reste du nuage, mélangeant quelques points C et D
-- vraisemblablement des cas atypiques dont les caractéristiques diffèrent
nettement du reste du jeu de données.

\begin{figure}[H]
    \centering
    \includegraphics[width=0.98\textwidth]{image/gradcam_approche2_resnet50.png}
    \caption{Grad-CAM sur l'Approche 2 (ResNet50, fine-tuning direct) :
    trois exemples par classe de densité, même protocole que la
    Figure~\ref{fig:gradcam}.}
    \label{fig:gradcam_approche2}
\end{figure}

Le résultat le plus frappant ici est \textbf{l'inverse} de ce qu'on
pourrait attendre. Les classes \textbf{A}, \textbf{B} et \textbf{D}
obtiennent des probabilités correctes élevées et stables sur leurs trois
exemples respectifs (A : 90.4\,\%, 89.2\,\%, 88.7\,\% ; B : 95.7\,\%,
87.9\,\%, 91.9\,\% ; D : 67.1\,\%, 92.1\,\%, 68.6\,\%), alors que la classe
\textbf{C} -- pourtant très majoritaire dans le jeu d'entraînement --
obtient des probabilités nettement plus faibles (17.4\,\%, 31.4\,\%,
17.0\,\%). Ce n'est pas une anomalie : c'est la confirmation visuelle
directe de ce qui a déjà été observé statistiquement pour cette approche
(Section~\ref{sec:approche2}) -- le modèle \og{}sur-route\fg{} une partie
des cas vers A, B et D au détriment de la confiance sur C, ce qui explique
à la fois son excellent rappel sur la classe A (0.95) et son rappel plus
faible sur la classe C (0.70) par rapport aux architectures qui, à
l'inverse, écrasent la classe A pour maximiser la précision sur C. Les
zones d'activation restent, comme pour l'Approche 3, concentrées sur le
tissu mammaire (autour du mamelon), sans réaction marquée au fond noir,
pour les quatre classes.

% ============================================================================
\chapter{Bugs identifiés et corrigés (récapitulatif méthodologique)}
\label{sec:bugs}

\begin{table}[h]
\centering
\small
\caption{Bugs identifiés lors de la relecture du code et leur correction}
\begin{tabular}{p{5cm}p{9cm}}
\toprule
Bug & Impact / correction \\
\midrule
Normalisation GLCM (débordement uint8) & Texture calculée sur des valeurs corrompues ; corrigé dans les 2 copies concernées \\
\texttt{FineTunedFeatureExtractor} (CNN) & Retournait des features spatiales au lieu des logits, faisait planter la loss \\
Filtrage vue manquant (\texttt{test\_hybrid\_model.py}) & Évaluation sur CC+MLO mélangés au lieu de la vue d'entraînement, score dilué \\
Appariement CC/MLO factice (test double-branche) & Les deux branches recevaient la même image dupliquée \\
Gel de branche inopérant (\texttt{hasattr} sur mauvais nom) & Le gel de branche ne s'appliquait jamais \\
Chargement de poids conditionnel à un seul \texttt{--backbone} & Une seule branche sur deux chargeait ses poids pré-entraînés \\
Chemin CONFIG obsolète (\texttt{training.py}) & Pointait vers un chemin Mac local au lieu du cluster \\
Résolution des chemins CBIS-DDSM (UUID) & Noms de fichiers réels ne correspondant pas aux CSV officiels, résolu par heuristique taille de fichier \\
Chemins CONFIG pointant vers un répertoire personnel d'un autre compte du cluster (\texttt{hybrid~train~mlo~rexnet/}) & Données d'entraînement et checkpoint de branche chargés depuis un compte externe au projet ; résultat correspondant (Modèle 3, Section~\ref{sec:modele3-invalide}) invalidé et retiré, code corrigé, puis modèle réentraîné avec succès (Test Acc 84\,\%, Focal Loss) sur les données propres au projet \\
\bottomrule
\end{tabular}
\end{table}

% ============================================================================
\chapter{Suivi opérationnel -- prochaines étapes}
\todo{ce chapitre est un outil de suivi interne, à mettre à jour au fur et à
mesure -- il n'est pas destiné à figurer tel quel dans la version finale du
mémoire. Une fois tous les entraînements terminés, le remplacer par un
véritable chapitre de conclusion et perspectives, synthétisant les résultats
obtenus approche par approche (Chapitre~\ref{ch:approches}), la comparaison
avec l'état de l'art (Chapitre~\ref{sec:etat-art-positionnement}), et les
limites/pistes d'amélioration déjà
identifiées dans ce document (classe A minoritaire, client drift du fédéré,
modèle 3 invalidé, etc.)}
\begin{itemize}
  \item Terminer modèle 2 GLCM (ResNet50+MLO)
  \item Récupérer la matrice de confusion de validation du modèle 1 GLCM (script en cours)
  \item Terminer Approche 2 ResNet50, Approche 6 (les 4 variantes)
  \item Rapport de classification complet ViT (Approche 2)
  \item \sout{Modèle 3 GLCM (RexNet150+Histogramme/MLO) -- invalidé, à réentraîner}
    -- fait : réentraîné avec Focal Loss, Test Acc 84\,\% (Section~\ref{sec:modele3-invalide})
  \item Une fois toutes les approches ci-dessus terminées : libérer un GPU et
    lancer l'entraînement fédéré (federated + centralized)
\end{itemize}

% ============================================================================
\begin{thebibliography}{9}

\bibitem{ref1}
Société canadienne du cancer,
\textit{Densité mammaire}.
Disponible sur : \url{https://cancer.ca/fr/treatments/tests-and-procedures/mammography/breast-density}

\bibitem{ref2}
Thelma-Radiologie,
\textit{Un outil de deep learning pour améliorer la classification des seins denses}.
Disponible sur : \url{https://www.thema-radiologie.fr/actualites/3229/un-outil-de-deep-learning-pour-ameliorer-la-classification-des-seins-denses.html}

\bibitem{ref3}
K. C. Ho et al.,
``Development and Validation of an AI-driven Mammographic Breast Density Classification Tool Based on Radiologist Consensus,''
\textit{Radiology: Artificial Intelligence}, vol. 3, no. 6, 2021.

\bibitem{ref4}
N. F. Boyd et al.,
``Mammographic Density and the Risk and Detection of Breast Cancer,''
\textit{New England Journal of Medicine}, vol. 356, no. 3, pp. 227--236, 2007.

\bibitem{ref5}
W. Eziddin, \textit{Segmentation itérative d'images par propagation de
connaissances dans le domaine possibiliste : application à la détection de
tumeurs en imagerie mammographique}.
Thèse de doctorat, Mines Telecom, 2012.

\bibitem{ref6}
S. R. Wellings,
``A Hypothesis of the Origin of Human Breast Cancer from the Terminal Ductal Lobular Unit,''
\textit{Pathol Res Pract}, vol. 166, pp. 515--535, 1980.
\url{https://doi.org/10.1016/S0344-0338(80)80248-2}

\bibitem{ref7}
S. M. McKinney \textit{et al.},
``International evaluation of an AI system for breast cancer screening,''
\textit{Nature}, vol. 577, pp. 89--94, 2020.

\bibitem{ref8}
L. Shen \textit{et al.},
``Deep Learning to Improve Breast Cancer Detection on Screening Mammography,''
\textit{Scientific Reports}, vol. 9, article 12495, 2019.
Disponible sur : \url{https://www.nature.com/articles/s41598-019-48995-4}

\bibitem{ref9}
H. F. R. Pereira \textit{et al.},
``Classification of Breast Density in Mammograms Using VGG16, ResNet50 and DenseNet,''
\textit{CIARP}, 2021.

\bibitem{ref10}
T. C. N. Alves \textit{et al.},
``Breast density classification using histograms,''
\textit{Springer}, 2014.

\bibitem{ref11}
R. Gatta \textit{et al.},
``Classification of breast density using Fuzzy c-means,''
\textit{Comput.\ Methods Programs Biomed.}, 2016.

\bibitem{ref12}
J. E. Lee \textit{et al.},
``Topographic segmentation and kNN modeling,''
\textit{Medical Physics}, vol. 44, no. 10, pp. 5231--5242, 2017.

\bibitem{ref13}
C. D. Lehman \textit{et al.},
``Mammographic Breast Density Assessment Using Deep Learning: Clinical Implementation,''
\textit{Radiology}, vol. 290, no. 1, pp. 52--58, 2019.
\todo{l'URL précédemment associée à cette clé (10.1148/ryai.210199)
appartient en réalité à Trace4BDensity (voir entrée séparée ci-dessous), pas
à cet article -- remplacer par le vrai DOI de Lehman et al.\ 2019 avant
remise finale.}

\bibitem{trace4bdensity}
S. Papa, F. Sardanelli, C. Salvatore \textit{et al.},
``Trace4BDensity: an AI tool for mammographic breast density classification,''
\textit{Radiology: Artificial Intelligence}, DOI: 10.1148/ryai.210199, 2022.
\todo{titre exact et année à confirmer sur la page de l'éditeur -- entrée
séparée créée suite à la correction du conflit de citation avec ref13
détecté lors de la relecture de la bibliographie}

\bibitem{ref14}
K. J. Geras \textit{et al.},
``High-Resolution Breast Cancer Screening with Multi-View Deep Convolutional Neural Networks,''
\textit{arXiv preprint arXiv:1711.05775}, 2017.

\bibitem{ref15}
H. Soltani,
``Détection du Cancer du Sein et Classification par une Approche intelligente,''
Thèse de Doctorat, Laboratoire LAMIS, Université Echahid Cheikh Larbi Tébessi de Tébessa, 2025.
Disponible sur : \url{https://lab-lamis-univ-tebessa.dz/Theses/Computer/Hama%20Soltani.pdf}

\bibitem{ref16}
\textit{Fusion collaborative de classifieurs : application à la détection de pathologies en mammographie}.
Disponible sur : \url{https://theses.hal.science/tel-03510707v1}

\bibitem{ref17}
``A classification framework for prediction of breast density using an ensemble of neural network classifiers.''
Disponible sur : \url{https://www.sciencedirect.com/science/article/pii/S0208521616302170}

\bibitem{ref18}
``Fully Automated Breast Density Segmentation and Classification Using Deep Learning,''
\textit{Diagnostics}, 2020.
Disponible sur : \url{https://www.mdpi.com/2075-4418/10/11/988}

\bibitem{ref19}
A. G. M. Al Mansour, F. Alshomrani, A. Alfahaid, A. T. M. Almutairi,
``MammoViT: A Custom Vision Transformer Architecture for Accurate BIRADS Classification in Mammogram Analysis,''
\textit{Diagnostics}, vol. 15, no. 3, article 285, 2025.
DOI: 10.3390/diagnostics15030285.
Disponible sur : \url{https://www.mdpi.com/2075-4418/15/3/285}

\bibitem{ref20}
O. Tanimola, O. Shobayo, O. Popoola, \& O. Okoyeigbo,
``Breast Cancer Classification Using Fine-Tuned SWIN Transformer Model on Mammographic Images,''
\textit{Analytics}, vol. 3, no. 4, pp. 461--475, 2024.
Disponible sur : \url{https://www.mdpi.com/2813-2203/3/4/26}

\bibitem{ref21}
L. W. Bassett, M. C. Mahoney, S. Apple, \& C. D'Orsi,
\textit{Breast Imaging: Expert Radiology Series E-Book}.
Elsevier Health Sciences, 2010.
Disponible sur : \url{https://www.uclahealth.org/departments/radiology/education/breast-imaging-teaching-resources/screening-mammogram/positioning-and-technique}

\bibitem{ref22}
I. Kumar, H. S. Bhadauria, J. Virmani, and S. Thakur,
``A classification framework for prediction of breast density using an ensemble of neural network classifiers,''
\textit{Biocybernetics and Biomedical Engineering}, vol. 37, no. 1, pp. 217--228, 2017.
Disponible sur : \url{https://www.researchgate.net/publication/313023655_A_classification_framework_for_prediction_of_breast_density_using_an_ensemble_of_neural_network_classifiers}

\bibitem{ref23}
M. F. S. Ho \textit{et al.},
``Density Map and Fuzzy Classification for BI-RADS Breast Tissue Composition,''
\textit{Computer Methods and Programs in Biomedicine}, vol. 199, article 105825, 2021.
Disponible sur : \url{https://doi.org/10.1016/j.cmpb.2020.105825}

\bibitem{ref24}
K. Jeffers \textit{et al.},
``Breast Cancer Risk and Mammographic Density Assessed with Semiautomated and Fully Automated Methods and BI-RADS,''
\textit{Radiology}, vol. 282, no. 2, pp. 348--355, 2017.
Disponible sur : \url{https://pubs.rsna.org/doi/10.1148/radiol.2016152062}

\bibitem{ref25}
L. Shen, L. Margolies, J. H. Rothstein, E. Fluder, R. McBride, and W. Zhang,
``End-to-End Training for Whole Image Breast Cancer Diagnosis Using an All Convolutional Design,''
\textit{arXiv preprint arXiv:1711.05775}, 2017.
Disponible sur : \url{https://arxiv.org/pdf/1711.05775}

\bibitem{ref26}
B. Boyer \textit{et al.},
``Automated 2D Breast Density Assessment: Correlation with Radiologist Consensus Visual Classification,''
\textit{Academic Radiology}, vol. 20, no. 12, pp. 1548--1554, 2013.

\bibitem{ref27}
C. Vachon \textit{et al.},
``Accuracy of BI-RADS Breast Density Category Definitions,''
\textit{Academic Radiology}, vol. 13, no. 10, pp. 1234--1241, 2006.
Disponible sur : \url{https://doi.org/10.1016/j.acra.2006.06.005}

\bibitem{ref28}
M. F. S. Ho \textit{et al.},
``Density Map and Fuzzy Classification for BI-RADS Breast Tissue Composition,''
\textit{Computer Methods and Programs in Biomedicine}, vol. 199, article 105825, 2021.
Disponible sur : \url{https://doi.org/10.1016/j.cmpb.2020.105825}
\todo{doublon exact de ref23 (même article) -- garder une seule des deux
clés et corriger les \texttt{\textbackslash cite} correspondants avant
remise finale}

\bibitem{ref29}
A. O. G. Guerriero \textit{et al.},
``Evaluation of Breast Density Assessment: A Comparative Study,''
\textit{The Breast}, vol. 16, no. 5, pp. 488--496, 2007.

\bibitem{ref30}
J. A. G. M. van Gils \textit{et al.},
``Reproducibility of BI-RADS Breast Density Categories,''
\textit{The Breast}, vol. 14, no. 3, pp. 210--216, 2005.

\bibitem{ref31}
P. M. Lauritzen \textit{et al.},
``AI in Mammography Screening: Workload Reduction and Screening Specificity,''
\textit{Radiology}, vol. 303, no. 2, pp. 281--289, 2022.
Disponible sur : \url{https://doi.org/10.1148/radiol.210948}

\bibitem{ref32}
A. Yala \textit{et al.},
``Using Deep Learning to Triage Screening Mammograms to Improve Radiologist Efficiency,''
\textit{Radiology}, vol. 293, no. 1, pp. 30--37, 2019.
Disponible sur : \url{https://doi.org/10.1148/radiol.2019182908}

\bibitem{ref33}
A. Rodríguez-Ruiz \textit{et al.},
``Detection of Breast Cancer with Mammography: Effect of an Artificial Intelligence Support System,''
\textit{Radiology}, vol. 290, no. 2, pp. 305--314, 2019.
Disponible sur : \url{https://doi.org/10.1148/radiol.2018181371}

% --- Références pour l'apprentissage fédéré et l'interprétabilité ---
\bibitem{ref_fedavg}
H. B. McMahan, E. Moore, D. Ramage, S. Hampson, and B. Agüera y Arcas,
``Communication-Efficient Learning of Deep Networks from Decentralized Data,''
\textit{Proceedings of the 20th International Conference on Artificial Intelligence and Statistics (AISTATS)}, PMLR 54:1273--1282, 2017.
Disponible sur : \url{https://arxiv.org/abs/1602.05629}

\bibitem{ref_haralick}
R. M. Haralick, K. Shanmugam, and I. Dinstein,
``Textural Features for Image Classification,''
\textit{IEEE Transactions on Systems, Man, and Cybernetics}, vol. SMC-3, no. 6, pp. 610--621, 1973.
Disponible sur : \url{https://doi.org/10.1109/TSMC.1973.4309314}

\bibitem{ref_gradcam}
R. R. Selvaraju, M. Cogswell, A. Das, R. Vedantam, D. Parikh, and D. Batra,
``Grad-CAM: Visual Explanations from Deep Networks via Gradient-Based Localization,''
\textit{Proceedings of the IEEE International Conference on Computer Vision (ICCV)}, pp. 618--626, 2017.
Disponible sur : \url{https://arxiv.org/abs/1610.02391}

\bibitem{ref_tsne}
L. van der Maaten and G. Hinton,
``Visualizing Data using t-SNE,''
\textit{Journal of Machine Learning Research}, vol. 9, pp. 2579--2605, 2008.

\bibitem{ref_gradcam_mammo1}
\todo{auteurs exacts à vérifier sur la page de l'éditeur}
``Explainable Mammogram Analysis with EfficientNetV2 and Grad-CAM++ for Robust Cancer Diagnosis,''
\textit{Diagnostics}, vol. 16, no. 1, article 105, 2025.
Disponible sur : \url{https://www.mdpi.com/2075-4418/16/1/105}

\bibitem{ref_gradcam_mammo2}
F. M. Talaat, S. A. Gamel, R. M. El-Balka, M. Shehata, and H. ZainEldin,
``Grad-CAM Enabled Breast Cancer Classification with a 3D Inception-ResNet V2: Empowering Radiologists with Explainable Insights,''
\textit{Cancers}, vol. 16, no. 21, article 3668, 2024.
DOI: 10.3390/cancers16213668.

\bibitem{ref_dietterich}
T. G. Dietterich,
``Ensemble Methods in Machine Learning,''
in \textit{Multiple Classifier Systems}, Lecture Notes in Computer Science, vol. 1857, pp. 1--15, Springer, 2000.
Disponible sur : \url{https://doi.org/10.1007/3-540-45014-9_1}

\bibitem{ref_focal_loss}
T.-Y. Lin, P. Goyal, R. Girshick, K. He, and P. Dollár,
``Focal Loss for Dense Object Detection,''
\textit{Proceedings of the IEEE International Conference on Computer Vision (ICCV)}, pp. 2980--2988, 2017.
Disponible sur : \url{https://arxiv.org/abs/1708.02002}

\bibitem{ref_rieke_fl_digital_health}
N. Rieke, J. Hancox, W. Li \textit{et al.},
``The Future of Digital Health with Federated Learning,''
\textit{npj Digital Medicine}, vol. 3, article 119, 2020.
Disponible sur : \url{https://www.nature.com/articles/s41746-020-00323-1}

\bibitem{ref_sheller_fl_brain}
M. J. Sheller, B. Edwards, G. A. Reina, J. Martin, S. Pati, A. Kotrotsou, M. Milchenko, W. Xu, D. Marcus, R. R. Colen, and S. Bakas,
``Federated Learning in Medicine: Facilitating Multi-Institutional Collaborations Without Sharing Patient Data,''
\textit{Scientific Reports}, vol. 10, article 12598, 2020.
Disponible sur : \url{https://www.nature.com/articles/s41598-020-69250-1}

\bibitem{ref_kaissis_fl_privacy}
G. A. Kaissis, M. R. Makowski, D. Rückert, and R. F. Braren,
``Secure, Privacy-Preserving and Federated Machine Learning in Medical Imaging,''
\textit{Nature Machine Intelligence}, vol. 2, pp. 305--311, 2020.
Disponible sur : \url{https://www.nature.com/articles/s42256-020-0186-1}

\bibitem{ref_roth_fl_breast_density}
H. R. Roth, K. Chang, P. Singh \textit{et al.},
``Federated Learning for Breast Density Classification: A Real-World Implementation,''
in \textit{Domain Adaptation and Representation Transfer, and Distributed and Collaborative Learning (DART/DCL Workshop, MICCAI)}, Lecture Notes in Computer Science, Springer, 2020.
Disponible sur : \url{https://arxiv.org/abs/2009.01871}

\end{thebibliography}

\end{document}
